% SPDX-FileCopyrightText: 2024-2026 Fabian Januszewski
% SPDX-License-Identifier: CC-BY-4.0

%======================================================================
% phenomena_gallery.tex -- [PG] the phenomena gallery: arithmetic phenomena of H^q(SL_2(Z); Sym^n) and of its level-N analogues,
% over F_ell, Z/ell^K and Z, arranged by mathematical phenomenon; each one shown by a short OC-script snippet with its real
% (generated) output, or by the project's runs.
%
%   Version 1.5.4 (2026-11-09): hacs instead of cohomology_app; BUILDING.md instead of CLAUDE.md (release audit S4).
%   Version 1.5.3 (slot 26, 2026-11-09): machine names and private paths removed (public release); no content change.
%   Version 1.5.2 (slot 26, 2026-11-09): confirm-read passed; two one-sentence repairs (sign check of ES_Z, Eis_1 twisted classes at N=49,64).
%   Version 1.5.1 (slot 26, 2026-11-09): referee fixes of sec:pg-nonsplit (split_on_P, sigma, coefficients of b, counts); tool modes sigma/extra/law.
%   Version 1.5.0 (slot 26, PG-NONSPLIT, 2026-11-08): new sec:pg-nonsplit, non-split torsion over Z (torsion-by-free and torsion-by-torsion), with a Galois reading; new snippet nonsplit_torsion_16_2; tool docs/latex/tools/pg_nonsplit_tors.py. NOT yet refereed.
%   Version 1.4.1 (slot 25, PG141, 2026-11-06): referee fixes R1--R6 of PG140_referee.md.
%   Version 1.4.0 (slot 25, PG140, 2026-11-06): the two l-adic objects (full localisation vs free quotient) named apart; census
%   edition 4.0 with section status, every census number recomputed by docs/latex/tools/pg_census_check.py.
%   Version 1.3.2 (slot 24, CELL4374-CLOSE, 2026-10-30): the (4374,3) cross-check landed. Version 1.3.1 (slot 24, DOCS-FIX, 2026-10-29): referee fixes P1--P9. Version 1.3.0 (slot 24, DOCS-PG). Reorganised by phenomenon (owner direction of 2026-10-29: logical, not
%   chronological; "think like a teacher"). The version history is Appendix app:pg-history.
%
%   STRUCTURE NOTE (1.2.1 -> 1.3.0). 1.2.1 had Part I "Console snippets" in the order the snippets were written (the ten of
%   slot 20, then one per verb of slots 21--22) and Part II "Phenomena from the runs" (slots 22--23). 1.3.0 has six parts, one
%   per phenomenon; each part shows its snippets first and then the run results on the same phenomenon:
%     I   Cohomology and its dimensions          sec:pg-dims (NEW), sec:pg-forms
%     II  Eisenstein and cuspidal classes        sec:pg-sign
%     III Congruences modulo ell                 sec:pg-ramanujan, -kummer, -kummerunit, -visible, -hidden, -dihedral, -mult2,
%                                                -theta, -loewy4; runs: -census34, -dihedral3, -satdef
%     IV  Level Gamma_0(N) and characters        sec:pg-leveln, -nebentypus; run: -mazur23
%     V   Hida theory: ordinary families         sec:pg-hida, -pow, -hidatower, -hidalambda; runs: -ord547, -hi11
%     VI  Integral cohomology and torsion        sec:pg-descent-torsion (NEW), -torsexp, -l23, -descent-congruence (NEW)
%     Appendix                                   sec:pg-missing, app:pg-history
%   Sections were renamed by phenomenon (a verb's name is now in parentheses); every \label of 1.2.1 is kept.
%
%   STATUS: NOT BINDING. This document defines nothing. Every definition it uses is a label of the binding
%   contract [C] (docs/latex/gl2_contract.tex), cited as {\ttfamily[C,~<label>]} (scanned by
%   docs/matrix-engine/tools/label_census.py), or of [E] hecke_extensions.tex, [U] hida_hecke_algebra.tex,
%   [EX] hecke_examples.tex, [L37] l37_loewy4_example.tex, [LG] loewy_and_galois_cohomology.tex, cited as
%   [E,~<label>] etc. (not scanned; each checked by hand at creation).
%
%   DRIFT GUARD. Every snippet is \VerbatimInput from notebook/examples/<name>.oc, every expected output
%   from docs/latex/phenomena_gallery/<name>.out, every runtime / RSS from phenomena_gallery/costs.tex.
%   The .out files and costs.tex are GENERATED by docs/latex/tools/gen_phenomena_outputs.sh <binary>;
%   `gen_phenomena_outputs.sh <binary> --check` re-runs every snippet and diffs against the committed .out.
%   Build: docs/latex/build_phenomena_gallery.sh <outdir>.
%======================================================================
\documentclass[11pt]{article}
\usepackage[T1]{fontenc}
\usepackage{lmodern}
\usepackage{microtype}
\usepackage[a4paper,margin=2.3cm]{geometry}
\usepackage{amsmath,amssymb}
\usepackage{amsthm}
\usepackage{booktabs}
\usepackage{array}
\usepackage{enumitem}
\usepackage{xcolor}
\usepackage{fvextra}
\usepackage[colorlinks=true,linkcolor=blue,citecolor=blue,urlcolor=blue]{hyperref}
\usepackage{xurl}

\setlength{\emergencystretch}{3em}
\newcommand{\Z}{\mathbb{Z}}
\newcommand{\Q}{\mathbb{Q}}
\newcommand{\F}{\mathbb{F}}
\newcommand{\Fl}{\mathbb{F}_\ell}
\newcommand{\Zl}{\mathbb{Z}_\ell}
\newcommand{\T}{\mathbb{T}}
\newcommand{\m}{\mathfrak{m}}
\newcommand{\Sym}{\operatorname{Sym}}
\newcommand{\SL}{\mathrm{SL}}
\newcommand{\GL}{\mathrm{GL}}
\newcommand{\Ga}{\Gamma}
\newcommand{\rhobar}{\bar\rho}
\newcommand{\code}[1]{\texttt{#1}}
\newcommand{\file}[1]{\texttt{#1}}
\newcommand{\pfile}[1]{\nolinkurl{#1}}
\newcommand{\PROVED}{\textsf{\small[PROVED]}}
\newcommand{\LIT}{\textsf{\small[LIT]}}
\newcommand{\CMPD}{\textsf{\small[CMPD]}}
\newcommand{\EMP}{\textsf{\small[EMP]}}
\newcommand{\PGCOND}{\textsf{\small[COND]}}
\newcommand{\PGSPEC}{\textsf{\small[SPEC]}}
\newtheorem*{proposition*}{Proposition}
\newtheorem*{theorem*}{Theorem}
\newcommand{\Hom}{\operatorname{Hom}}
\newcommand{\smat}[4]{\bigl(\begin{smallmatrix}#1&#2\\#3&#4\end{smallmatrix}\bigr)}
\newcommand{\pgtodo}[1]{\textcolor{red}{\textsf{\small[TODO: #1]}}}

% the generated costs: \pgcost{<snippet>}
\makeatletter
% GENERATED by docs/latex/tools/gen_phenomena_outputs.sh -- do not edit by hand.
% binary: 21a00a018b81a7c2...  host: x86_64  date: 2026-11-08
\expandafter\def\csname pgcost@cohomology_dimensions\endcsname{0.16\,s, 26\,MiB}
\expandafter\def\csname pgcost@forms_vs_cohomology_24_7\endcsname{0.01\,s, 25\,MiB}
\expandafter\def\csname pgcost@eisenstein_sign_cone_10_691\endcsname{0.36\,s, 49\,MiB}
\expandafter\def\csname pgcost@ramanujan_10_691\endcsname{0.38\,s, 58\,MiB}
\expandafter\def\csname pgcost@kummer_ribet_30_37\endcsname{0.31\,s, 32\,MiB}
\expandafter\def\csname pgcost@kummer_unit_37_59\endcsname{0.03\,s, 22\,MiB}
\expandafter\def\csname pgcost@visible_congruence_48_7\endcsname{0.81\,s, 71\,MiB}
\expandafter\def\csname pgcost@hidden_congruence_300_23\endcsname{4.63\,s, 370\,MiB}
\expandafter\def\csname pgcost@dihedral_wilton_10_23\endcsname{3.14\,s, 137\,MiB}
\expandafter\def\csname pgcost@multiplicity_two_34_5\endcsname{0.55\,s, 51\,MiB}
\expandafter\def\csname pgcost@dickson_theta_16_7\endcsname{0.13\,s, 35\,MiB}
\expandafter\def\csname pgcost@loewy4_258_37\endcsname{4.25\,s, 725\,MiB}
\expandafter\def\csname pgcost@level_n_old_forms_2_7\endcsname{0.00\,s, 21\,MiB}
\expandafter\def\csname pgcost@nebentypus_7_13\endcsname{0.00\,s, 21\,MiB}
\expandafter\def\csname pgcost@hida_control_13\endcsname{2.39\,s, 278\,MiB}
\expandafter\def\csname pgcost@pow_ordinary_rank_13\endcsname{0.14\,s, 46\,MiB}
\expandafter\def\csname pgcost@hida_tower_13_10\endcsname{0.02\,s, 24\,MiB}
\expandafter\def\csname pgcost@hida_lambda_11_12\endcsname{0.00\,s, 21\,MiB}
\expandafter\def\csname pgcost@descent_torsion_beta\endcsname{1.94\,s, 543\,MiB}
\expandafter\def\csname pgcost@descent_congruence_primes\endcsname{0.17\,s, 31\,MiB}
\expandafter\def\csname pgcost@ramanujan_691_over_Z\endcsname{0.35\,s, 34\,MiB}
\expandafter\def\csname pgcost@nonsplit_torsion_16_2\endcsname{0.67\,s, 152\,MiB}

\newcommand{\pgcost}[1]{\@ifundefined{pgcost@#1}{(not measured)}{\csname pgcost@#1\endcsname}}
\makeatother

% a snippet: its source, then its recorded output
\newcommand{\snippet}[1]{%
  \par\smallskip\noindent\textbf{The snippet} \file{notebook/examples/\detokenize{#1}.oc}:\par
  \VerbatimInput[breaklines=true,breakanywhere=true,fontsize=\footnotesize,frame=single,framesep=2mm,
    rulecolor=\color{black!40}]{../../notebook/examples/#1.oc}
  \par\noindent\textbf{Expected output} (generated file \file{phenomena\_gallery/\detokenize{#1}.out}):\par
  \VerbatimInput[breaklines=true,breakanywhere=true,fontsize=\footnotesize,frame=leftline,framesep=2mm,
    rulecolor=\color{blue!50}]{phenomena_gallery/#1.out}
  \par\noindent\textbf{Cost} (wall time, peak RSS; reference box, binary named in \file{costs.tex}): \pgcost{#1}.\par\medskip}

\newenvironment{lookfor}{\par\noindent\textbf{What to look for.}\begin{itemize}[leftmargin=1.6em,itemsep=1pt]}{\end{itemize}}

\title{Phenomena gallery: the arithmetic of $H^{*}(\SL_2(\Z);\Sym^{n})$ over $\Fl$, $\Z/\ell^{K}$ and $\Z$,\\ in twenty-one console snippets and the project's runs}
\author{Fabian Januszewski\\ Universit\"at Paderborn\\ \texttt{fabian.januszewski@math.uni-paderborn.de}}
\date{Version 1.5.4 --- 2026-11-09. Not binding. Version history: Appendix~\ref{app:pg-history}.}

\begin{document}
\maketitle

\begin{abstract}
This gallery shows, one phenomenon at a time, what the project computes about the cohomology of $\SL_2(\Z)$ (and of $\Ga_0(N)$) with
coefficients in the symmetric powers $\Sym^{n}$ of the standard representation, together with its Hecke operators. It is arranged by
mathematics, in the order in which one would teach it:
\begin{enumerate}[label=\Roman*.,leftmargin=2.2em,itemsep=1pt]
\item \emph{Cohomology and its dimensions}: $\dim H^{1}=2\dim S_k+1$ (Eichler--Shimura), and where small primes break the count;
\item \emph{Eisenstein and cuspidal classes}: the involution $\iota^{*}$ and the sign of the Eisenstein class, the boundary, the compactly
  supported cohomology;
\item \emph{Congruences modulo $\ell$}: Ramanujan's $691$ as an Eisenstein ideal, irregular primes and Herbrand--Ribet, visible and hidden
  congruences, a dihedral residual representation, multiplicity two, the $\theta$-operator, Loewy length~$4$ --- and what the congruence census found;
\item \emph{Level $\Ga_0(N)$ and characters}: old forms, nebentypus, Mazur's Eisenstein ideal at $p=2,3$;
\item \emph{Hida theory}: ordinary parts constant along a weight class, control modulo $\ell^{m}$, to $n=6910$ at $\ell=691$;
\item \emph{Integral cohomology and torsion}: $H^{1}$ over $\Z$ by the integral descent, its torsion primes, the exponent law $\beta(n,\ell)$
  (Callegaro--Cohen--Salvetti), checked to $n=10^{4}$, a first non-split extension of Hecke modules between torsion and free part, and the Eisenstein congruence primes $691$ and $5$ read from the Hecke operators
  over~$\Z$.
\end{enumerate}
Each phenomenon is stated in a few sentences, cited to the binding contract [C] (and [E], [U]) by label, shown by the exact command and its
\emph{real} output --- every listing and output is included from the files themselves and regenerated by a script, so the document and the
program cannot diverge --- and followed by what to look for in that output. Phenomena that need hours and tens of gigabytes are reported from
the runs, with their sources, after the snippets of the same part. Nothing here is new mathematics unless it says so, and each run result is
tagged as proved, certified computation, or empirical pattern; where a literature search found a source, it is cited.
\end{abstract}

\tableofcontents

\section{Introduction}\label{sec:pg-intro}

\paragraph{The objects.} For an even $n\ge0$ and a prime $\ell$ the console computes
$H^{q}=H^{q}(\SL_2(\Z);\Sym^{n}\Fl^{2})$, $q\in\{0,1,2\}$, on the seven-vertex orbifold model {\ttfamily[C,~def:model7]} with the
symmetric-power convention of {\ttfamily[C,~prop:symconv]}, together with the Hecke operators $T_p$ (and $U_\ell=T_\ell$ by the same
code path), defined by the double-coset chain map of {\ttfamily[C,~def:hecke]}, {\ttfamily[C,~eq:heckechain]}. The weight of the cell
$(n,\ell)$ is $k=n+2$. For $\ell\ge5$ and $0<n<\ell$ (and in general, for $n>0$, after the Eichler--Shimura isomorphism over $\Q$), $H^{1}$ has
dimension $2\dim S_k+1$: two classes for each cusp eigenform of level one and weight $k$, one in each eigenspace of the involution
$\iota^{*}$, and one Eisenstein class (\LIT). The eigenvalue of $T_p$ on the Eisenstein class is $1+p^{k-1}$; on the classes of a cusp
eigenform $f=\sum a_n q^n$ it is $a_p\bmod\ell$. All of this is visible in the snippets, and every congruence below is a statement about
the $\Fl$-algebra generated by the $T_p$ on $H^{1}$ (or its $\Z/\ell^{K}$-analogue), with the eigenforms in the background.
The same verbs work over $\Z/\ell^{K}$ (\code{set coefficient\_ring = "Z/ell\^{}2"}; Parts~\ref{part:pg-congruences} and~\ref{part:pg-hida}),
and the integral descent computes over $\Z$, at level one and at level $\Ga_0(N)$ (\code{set coefficient\_ring = "Z"}; Part~\ref{part:pg-integral}).

\paragraph{How the gallery is organised.} The six parts follow the mathematics, not the history of the code. Part~\ref{part:pg-cohomology}
computes the groups and their dimensions; Part~\ref{part:pg-eisenstein} splits $H^{1}$ into its Eisenstein and cuspidal pieces;
Part~\ref{part:pg-congruences} studies how the two pieces (and different cusp forms) are glued modulo $\ell$ --- the local Hecke algebras;
Part~\ref{part:pg-level} changes the group to $\Ga_0(N)$ and twists by characters; Part~\ref{part:pg-hida} moves the weight and follows
the ordinary part along $\ell$-adic families; Part~\ref{part:pg-integral} returns to $\Z$-coefficients, where the torsion of the integral
cohomology explains the jumps of Part~\ref{part:pg-cohomology} and the congruence primes of Part~\ref{part:pg-congruences} reappear in
the integral Hecke operators. Within a part, the snippets come first (seconds each) and the run results second
(\S\ref{sec:pg-runs}). Each section ends with a box \emph{What to look for}. Read in order, the parts form a short course; each section
can also be read alone, with its back references.

\paragraph{How to run a snippet.} With a built binary \file{hacs} (\file{BUILDING.md}; \file{cohomology\_app} is a link to it):
\begin{itemize}[leftmargin=1.6em]
\item \emph{as a script:} \code{hacs --script notebook/examples/ramanujan\_10\_691.oc} (exit code $0$; a failed
  \code{assert} exits $4$);
\item \emph{in the console:} \code{hacs -i}, then paste the lines, or \code{hacs --script <file> --console} to run
  the snippet and keep the session (every variable stays bound; \code{?x} describes a value, \code{help <verb>} documents a verb);
\item \emph{in the notebook:} configure with \code{-DOC\_BUILD\_NOTEBOOK=ON} and start the gateway
\begin{Verbatim}[fontsize=\small]
python3 notebook/gateway/oc_gateway.py --binary <build>/hacs
\end{Verbatim}
  (loopback by default; it prints a one-time token), open the page and paste a snippet into a cell. The notebook serves only verbs tagged \code{read\_only}
  (\file{Console/serve/read\_only\_verbs.tsv}) and refuses \code{include}, \code{quit} and \code{exit}; \code{set} of a mathematical
  setting is allowed. All twenty-one snippets satisfy these rules (none uses \code{include}; the Hida and descent snippets use \code{set coefficient\_ring}).
\end{itemize}
The language is documented in the CLI manual [CLI] (\file{docs/latex/cli\_manual.pdf}) and in the generated reference
\file{docs/console/console.md}: a handle such as \code{cohomology(...)} is lazy and computes on first access to a member; \code{print}
prints its arguments separated by spaces; \code{/} on integers is exact (an error on a remainder), so the snippets write $\lfloor e/2\rfloor$
as \code{(e - e \% 2) / 2}; integer overflow is an error, so powers modulo $\ell$ are taken by a loop.

\paragraph{The drift guard.} Each listing is read from \file{notebook/examples/<name>.oc} and each output from
\file{docs/latex/phenomena\_gallery/<name>.out}. The outputs and the cost line are written by
\file{docs/latex/tools/gen\_phenomena\_outputs.sh <binary>}, which runs every snippet under a $4$\,GiB RSS watchdog;
\code{gen\_phenomena\_outputs.sh <binary> --check} re-runs them and diffs against the committed outputs (costs are measured, never compared).
The PDF is built by \file{docs/latex/build\_phenomena\_gallery.sh <outdir>}.

\paragraph{Tags.} \PROVED{} proved in a cited document; \LIT{} classical literature, stated for orientation; \CMPD{} computed by the
snippet shown (and only that). No statement of this document is binding; [C] is.

\begin{center}\footnotesize
\begin{tabular}{@{}rlll@{}}
\toprule
\S & phenomenon & cell(s) $(n,\ell)$ & snippet \\
\midrule
\multicolumn{4}{@{}l}{\emph{I.\ Cohomology and its dimensions}}\\
\ref{sec:pg-dims} & $\dim H^{q}$, Eichler--Shimura, small $\ell$ & $n\le24$, $\ell=691,5$ & \file{cohomology\_dimensions} \\
\ref{sec:pg-forms} & cohomology against $q$-expansions & $(22,7)$ & \file{forms\_vs\_cohomology\_24\_7} \\
\multicolumn{4}{@{}l}{\emph{II.\ Eisenstein and cuspidal classes}}\\
\ref{sec:pg-sign} & $\iota^{*}$, ``Eisenstein $=+$'', boundary and cone & $(10,691)$ & \file{eisenstein\_sign\_cone\_10\_691} \\
\multicolumn{4}{@{}l}{\emph{III.\ Congruences modulo $\ell$}}\\
\ref{sec:pg-ramanujan} & Ramanujan's congruence, the Eisenstein ideal & $(10,691)$ & \file{ramanujan\_10\_691} \\
\ref{sec:pg-kummer} & irregular primes, $\zeta(1-k)$, Herbrand--Ribet & $(10,37)$, $(30,37)$ & \file{kummer\_ribet\_30\_37} \\
\ref{sec:pg-kummerunit} & Kummer's test on cyclotomic units & $\ell\in\{37,59\}$ & \file{kummer\_unit\_37\_59} \\
\ref{sec:pg-visible} & a visible congruence & $(48,7)$ & \file{visible\_congruence\_48\_7} \\
\ref{sec:pg-hidden} & a hidden congruence & $(300,23)$ & \file{hidden\_congruence\_300\_23} \\
\ref{sec:pg-dihedral} & a dihedral $\rhobar$: Wilton's congruence & $(10,23)$ & \file{dihedral\_wilton\_10\_23} \\
\ref{sec:pg-mult2} & multiplicity one fails mod $5$ & $(34,5)$ & \file{multiplicity\_two\_34\_5} \\
\ref{sec:pg-theta} & the $\theta$-operator and the Dickson sequence & $(16,7)\to(24,7)$ & \file{dickson\_theta\_16\_7} \\
\ref{sec:pg-loewy4} & Loewy length $4$ mod $37$ & $(258,37)$ & \file{loewy4\_258\_37} \\
\multicolumn{4}{@{}l}{\emph{IV.\ Level $\Ga_0(N)$ and characters}}\\
\ref{sec:pg-leveln} & old forms at level $2$ & $(10,7)$ & \file{level\_n\_old\_forms\_2\_7} \\
\ref{sec:pg-nebentypus} & a character, the parity obstruction & $(2,13)$ & \file{nebentypus\_7\_13} \\
\multicolumn{4}{@{}l}{\emph{V.\ Hida theory}}\\
\ref{sec:pg-hida} & Hida's control theorem mod $13^{2}$ & $(10{+}12j,13)$ & \file{hida\_control\_13} \\
\ref{sec:pg-pow} & the ordinary rank in one call & $(10{+}12j,13)$ & \file{pow\_ordinary\_rank\_13} \\
\ref{sec:pg-hidatower} & a Hida tower cell & $(10,13)$ & \file{hida\_tower\_13\_10} \\
\ref{sec:pg-hidalambda} & control on forms & $k=12$, $122$; $\ell=11$ & \file{hida\_lambda\_11\_12} \\
\multicolumn{4}{@{}l}{\emph{VI.\ Integral cohomology and torsion}}\\
\ref{sec:pg-descent-torsion} & $H^{1}$ over $\Z$, torsion primes, $\beta(n,\ell)$ & $n=4..60$, level $1$ & \file{descent\_torsion\_beta} \\
\ref{sec:pg-descent-congruence} & congruence primes, boundary pair & $n=10$; $\Ga_0(11)$, $n=0$ & \file{descent\_congruence\_primes} \\
\ref{sec:pg-nonsplit} & non-split torsion, Galois & $\Ga_0(8)$, $16$, $27$, $9$ & \file{nonsplit\_torsion\_16\_2} \\
\bottomrule
\end{tabular}
\end{center}

\section{Snippets and runs: two kinds of evidence}\label{sec:pg-runs}\label{part:pg-runs}

A snippet reproduces its phenomenon in seconds, on any build. Each part then records the phenomena of the same kind found in the project's \emph{runs}:
\begin{itemize}[leftmargin=1.6em,itemsep=1pt]
\item \textbf{census edition~4.0} (\pfile{docs/matrix-engine/results/congruence_census/edition3/4.0/}), in four sections of different status:
  \begin{itemize}[itemsep=0pt]
  \item the \emph{main} section (\pfile{rows40.tsv.gz}; $279\,093$ rows in $13\,505$ cells: the $27$ primes $5\le\ell\le103$, $\ell=547$, $\ell=691$,
    even $n\le1000$) and \emph{section b1} (\pfile{rows40_b1.tsv.gz}; $169\,885$ rows in $9500$ cells: the $19$ primes $107\le\ell\le199$, even
    $n\le1000$) are the edition \textbf{signed} by the owner on 2026-10-27: $448\,978$ rows (one per maximal ideal of a cell's Hecke algebra) in
    $23\,005$ cells, $46$ primes. Edition~3.4 (signed 2026-10-15, $n\le700$) is contained in the main section row for row; its saturation
    columns were extended by the step SATC of 3.5. In $580$ main cells (large $n$ at $\ell\ge73$, $547$, $691$) the record has no $U_\ell$ by design;
  \item \emph{section b2} (\pfile{rows40_b2.tsv.gz}, 4.0-pre.3; $152\,760$ rows in $11\,911$ cells, $35$ primes $211\le\ell\le419$) is a
    \textbf{partial pre-release} addendum, not signed: complete for $\ell\le347$, $411$ of $6000$ cells for $349\le\ell\le419$;
  \item the \emph{L23 section} (\pfile{README_l23.md}: the $\ell\in\{2,3\}$ cells, $1052$ cells, $n\le700$, and the DEEP-1 tables, $12\,220$
    cells, $n\le10^{4}$) is a \textbf{pre-release} addendum, not signed; its tables are byte-identical to those of 3.5;
  \end{itemize}
\item the \textbf{integral census, edition~1} (\pfile{docs/integral_census/results/ed1/}; signed by the owner 2026-11-03): $1500$ descent cells
  $(N,n)$ of $H^{q}(\Ga_0(N);\Sym^{n}\Z^{2})$ (level one with even $n\le400$, $2\le N\le100$ with even $n\le20$, and the boundary-pair grid);
  the gallery quotes no number from it, its integral examples are the snippets of Part~\ref{part:pg-integral};
\item the saturation runs, the second-cluster Hida run ACL2O, and the level-$N$ runs at $p=2,3$.
\end{itemize}
Every census number in the run sections says which section it comes from; ``the census'' without qualification means the signed edition
(main $+$ b1). The numbers are recomputed from the rows files by \file{docs/latex/tools/pg\_census\_check.py}, which reads the printed literals out of this text and compares each with the data (\code{PG\_CENSUS\_CHECK \dots{} verdict=OK}); ordinarity, which is not a census column, it reads from the cells' records.
These cells take minutes to days and tens of gigabytes, so the run sections have no snippets: every figure is quoted from a committed report or
data file, named in each section.

\paragraph{Tags in the run sections.} \CMPD{} there means \emph{computed and certified on exactly the cells stated}. One more tag: \EMP{} an empirical
pattern on the cells stated, not proved and possibly false outside them. The phenomenology catalogue
\pfile{docs/phenomenology/phenomena_catalogue.md} (entries PH-01 \dots{} PH-48) lists every phenomenon with its evidence, sources and status,
including those that are documented elsewhere. The run sections cover the ones that no LaTeX document described on 2026-10-19. The catalogue's
version~2 adds the literature verdict of every entry, and \pfile{docs/phenomenology/FINDINGS_SUMMARY.md} ranks the items for which no prior
source was found.

\paragraph{Notation.}
\begin{itemize}[leftmargin=1.6em,itemsep=1pt]
\item $\Ga=\SL_2(\Z)$, $V_n=\Sym^{n}\Z^{2}$, $M_n=V_n\otimes\Fl$, weight $k=n+2$. A cell is written $(n,\ell)$, degree first; level-$N$ cells
  are written $(N,p)$ and say so.
\item A \emph{block} is a generalised eigenspace $H^{1}(\Ga;M_n)_\m$, of dimension $D$. $A=A_\m$ is the image of the Hecke algebra in its
  endomorphisms (the \emph{census algebra}).
\end{itemize}

\paragraph{Two $\ell$-adic objects: the full localisation and its free quotient.} Over $\Zl$ the cohomology has torsion
(Part~\ref{part:pg-integral}), so there are two lattices to localise, and the gallery keeps them apart:
\begin{itemize}[leftmargin=1.6em,itemsep=1pt]
\item $\mathbb L_\m:=H^{1}(\Ga;V_n\otimes\Zl)_\m$, the \emph{full} localisation, torsion included: as a $\Zl$-module
  $\mathbb L_\m\cong\Zl^{f}\oplus\bigoplus_i\Z/\ell^{e_i}$. For $\ell\ge5$ its reductions are the cohomology with finite coefficients,
  $\mathbb L_\m/\ell^{K}=H^{1}(\Ga;V_n\otimes\Z/\ell^{K})_\m$ and in particular $\mathbb L_\m/\ell=H^{1}(\Ga;M_n)_\m$, the block
  ({\ttfamily[C,~lem:freepart]}, with {\ttfamily[C,~cl:bockstein]}: $H^{2}$ has no $\ell$-torsion). Its torsion is non-zero only on the Dickson locus
  (\S\ref{sec:pg-descent-torsion}).
\item $\Lambda^{\mathrm{fr}}_\m:=\mathbb L_\m/\mathrm{tors}\cong\Zl^{f}$, its \emph{free quotient} ($\mathbb L'_\m$ of {\ttfamily[C,~sec:satlevelone]}). On it,
  $\T_\m$ is the $\Zl$-algebra generated by the $T_p$ (the image of the Hecke algebra in $\mathrm{End}(\Lambda^{\mathrm{fr}}_\m)$,
  [C]'s $\T^{\mathbb L}_\m$), $\T_\m^{\mathrm{sat}}$ its saturation and $\mathsf A_K$ its image in $\mathrm{End}(\Lambda^{\mathrm{fr}}_\m/\ell^{K})$
  {\ttfamily[C,~def:satdata]}, which is stated for a \emph{free} lattice $L$; $\delta=\mathrm{length}(\T^{\mathrm{sat}}_\m/\T_\m)$ is the
  saturation defect {\ttfamily[C,~prop:satdefect]}.
\end{itemize}
When the block has no torsion ($\mathbb L_\m=\Lambda^{\mathrm{fr}}_\m$, the generic case) the two coincide, $A=\mathsf A_1$, and every statement
below can be read on either. When it has torsion they differ, and the statistics live on different objects:
\begin{itemize}[leftmargin=1.6em,itemsep=1pt]
\item the census algebra $A$, its Hilbert function, Loewy length and module structure live on the \emph{full} mod-$\ell$ block
  $\mathbb L_\m/\ell=H^{1}(\Ga;M_n)_\m$ (the extension record of the cell, \file{HeckeExt/HeckeExtDriver.cpp});
\item the $\Z/\ell^{K}$ module route (\code{hecke\_ring\_algebra}, \S\ref{sec:pg-visible}) runs on the full $\mathbb L_\m/\ell^{K}$, torsion
  summands included, reads $A$ there (the console's \code{A1}; it equals $\mathsf A_1$ only on a torsion-free block), and on a torsion block reads the free part separately, as $\ell^{t}(\mathbb L_\m/\ell^{K})
  \cong\Lambda^{\mathrm{fr}}_\m/\ell^{K-t}$ ($t$ the torsion exponent, {\ttfamily[C,~lem:freepart]}, the verdict \code{torsion\_free\_stable});
\item the saturation statistics $\delta$, $\dim\mathrm{Tor}_1(\T^{\mathrm{sat}}/\T,\Fl)$ and its Loewy length, and the census flag
  \code{hidden\_congruence}, live on $\Lambda^{\mathrm{fr}}_\m$ ($\delta$ and the saturation data as [C] defines them, the hidden flag as in
  [LG,~def:hiddencong]), and the census computes them \emph{only on torsion-free
  blocks}: a torsion block is refused by name (\code{not\_computed:torsion\_block}, \code{not\_stable:torsion\_block}), so no census value
  mixes the two objects;
\item (Concordance of notation: [E] and [LG] write $\Lambda_\m$ for the full lattice, our $\mathbb L_\m$; [C] writes $\mathbb L'_\m$ for our $\Lambda^{\mathrm{fr}}_\m$.)
\item the integral descent and the integral census (Part~\ref{part:pg-integral}) compute the full group over $\Z$, with the Hecke operators
  on its free quotient (the integral characteristic polynomials) and on each $\ell$-primary torsion part.
\end{itemize}
Since $\mathbb L_\m\to\Lambda^{\mathrm{fr}}_\m$ is onto and Hecke-equivariant, $\Lambda^{\mathrm{fr}}_\m/\ell$ is a quotient of the block, so the
algebra $\mathsf A_1$ of the free part is a quotient of $A$: $\dim\mathsf A_1\le\dim A$, with equality on torsion-free blocks.

%======================================================================
\part{Cohomology and its dimensions}\label{part:pg-cohomology}\label{part:pg-snippets}

Before any arithmetic, the groups. The group $\Ga=\SL_2(\Z)$ is the amalgam $\Z/4\ast_{\Z/2}\Z/6$ {\ttfamily[C,~cl:amalgam]}; the project computes its
cohomology on a seven-vertex model of the orbifold $\Ga\backslash\mathfrak H$ {\ttfamily[C,~def:model7]}. This part shows the dimensions over a
field, the classical count that explains them, and the first sign that small primes behave differently.

%======================================================================
\section{The dimensions of \texorpdfstring{$H^{q}$}{H^q} and the Eichler--Shimura count}\label{sec:pg-dims}

The centre $\{\pm I\}$ of $\Ga$ acts on $\Sym^{n}$ by $(-1)^{n}$, so for $n$ odd and $2$ invertible all cohomology vanishes {\ttfamily[C,~rmk:negI]}.
For $n$ even and a field $F$ in which $6$ is invertible the finite stabilisers $\Z/4$, $\Z/6$ have no higher cohomology, so $H^{q}=0$ for $q\ge2$,
$H^{0}$ is the space of invariants ($0$ for $n>0$), and everything sits in $H^{1}$. Its size is the Eichler--Shimura theorem
{\ttfamily[C,~thm:classical]} (\LIT): for even $n>0$ and $k=n+2$,
\[
  H^{1}(\Ga;\Sym^{n}\mathbb C^{2})\;\cong\;S_k\oplus\overline{S_k}\oplus\mathbb C\cdot E_k,\qquad
  \dim H^{1}=2\dim S_k+1,
\]
\[
  \dim S_k=\lfloor k/12\rfloor-[k\equiv2\bmod12],
\]
two classes for each cusp form and one for the Eisenstein series. Over $\Fl$ with $\ell>n+1$ the integral cohomology has no $\ell$-torsion
(\S\ref{sec:pg-descent-torsion}), so the same count holds; part~(1) of the snippet checks it for $n\le24$ at $\ell=691$, computing
$\dim S_k$ from the formula, not from a table.

Part~(2) takes $\ell=5$. Now $\Z/4$ and $\Z/6$ still have no $5$-torsion, so $H^{2}$ stays $0$, but $H^{0}(\Ga;\Sym^{n}\F_5^{2})$ is non-zero
exactly when $n$ is a \emph{Dickson degree} at $5$, i.e.\ $n\in6\,\Z_{\ge0}+20\,\Z_{\ge0}$: the invariants of $\SL_2(\F_5)$ are generated by the
Dickson invariants of degrees $\ell+1=6$ and $\ell(\ell-1)=20$ {\ttfamily[C,~def:dicksoninv]}. By the universal coefficient theorem each such
invariant is the shadow of a cyclic summand $\Z/5^{a}$ of $H^{1}(\Ga;\Sym^{n}\Z^{2})$ {\ttfamily[C,~cl:bockstein]}, which raises $\dim H^{0}$
and $\dim H^{1}$ over $\F_5$ by one each: at $n=12$, $\dim H^{1}=2$ against the characteristic-zero count $1$. Part~\ref{part:pg-integral}
computes these torsion groups over $\Z$ directly.

\snippet{cohomology_dimensions}

\begin{lookfor}
\item odd $n$: \code{dim H\^{}q = [0, 0, 0]}; $n=0$: only $H^{0}$, the constants;
\item even $n>0$ at $\ell=691$: \code{dim H\^{}1 = 2 dim S\_k + 1}, e.g.\ $3$ at $k=12$ ($\Delta$), $5$ at $k=24$, $1$ at $k=14$ (no cusp form);
\item at $\ell=5$: \code{dim H\^{}0 = 1} exactly at $n=0,6,12,18,20,24$, and for $n>0$ among these $\dim H^{1}$ exceeds $2\dim S_k+1$ by one
  (at $n=0$, $H^{1}=0$).
\end{lookfor}

%======================================================================
\section{Cohomology against \texorpdfstring{$q$}{q}-expansions (\texttt{forms\_algebra})}\label{sec:pg-forms}

\code{forms\_algebra(k, ell)} computes $S_k$ over $\Fl$ from $q$-expansions (Miller basis) with the matrices of $T_p$. By Eichler--Shimura
{\ttfamily[C,~thm:classical]} (\S\ref{sec:pg-dims}) the characteristic polynomial of $T_p$ on $H^{1}(\SL_2(\Z);\Sym^{k-2}\Q^{2})$ is the square of the one on $S_k$ times $x-(1+p^{k-1})$;
it reduces to $\Fl$ when $H^{1}$ over $\Fl$ is the reduction of the lattice (no $\ell$-torsion in $H^{1}$ over $\Z$; here $\dim=5=2\cdot2+1$ shows it).
At $(k,\ell)=(24,7)$ the snippet checks this identity for $p=2$ as an \code{assert} (\CMPD): two independent routes, $q$-expansions and the cohomology of the
seven-vertex model, give the same polynomial.

\snippet{forms_vs_cohomology_24_7}

\begin{lookfor}
\item \code{d = 2}: $\dim S_{24}=2$; the cohomology has dimension $5=2\cdot2+1$;
\item the \code{assert} passes: the identity is checked, not displayed.
\end{lookfor}

%======================================================================
\part{Eisenstein and cuspidal classes}\label{part:pg-eisenstein}

$H^{1}$ contains one Eisenstein class and two classes per cusp form (\S\ref{sec:pg-dims}). This part tells them apart: by an involution,
and by the boundary of the Borel--Serre compactification, where only the Eisenstein class survives. The compactly supported cohomology
$H^{1}_{c}$ and the interior cohomology $H^{1}_{!}$, the cuspidal part, come out of the same long exact sequence.

%======================================================================
\section{The sign of the Eisenstein class, the boundary and the cone}\label{sec:pg-sign}

The involution of the upper half-plane $z\mapsto-\bar z$ induces an involution on $H^{q}$; the project's production involution is
$\iota^{*}:=-\iota^{*}_{\mathrm{geo}}$ {\ttfamily[C,~def:iota]}, and the sign was fixed so that the Eisenstein class lies in $H^{1}_{+}$:
a convention by definition, not a theorem {\ttfamily[C,~rmk:iotasign]}. For $2\in R^{\times}$, $H^{q}=H^{q}_{+}\oplus H^{q}_{-}$
{\ttfamily[C,~prop:iotasplit]}, and for $6\in R^{\times}$ both summands are Hecke-stable. Each cusp eigenform contributes one class to each
sign.

At $(10,691)$ the split is $\dim H^{1}_{+}=2$ (Eisenstein, $\Delta^{+}$) and $\dim H^{1}_{-}=1$ ($\Delta^{-}$). The non-split extension of Ramanujan's congruence
(\S\ref{sec:pg-ramanujan}, the first section of Part~\ref{part:pg-congruences}) lives \emph{entirely in $H^{1}_{+}$}: there the module is cyclic with Loewy series $[1,1]$, while $H^{1}_{-}$ is a
single semisimple line. This is the shape that Harder's and Ribet's constructions predict for an Eisenstein congruence
[E,~cl:harder].

The boundary of the Borel--Serre compactification is computed by the circle model {\ttfamily[C,~prop:circlemodel]}; $T_p$ acts on $H^{*}_{\partial}$
through the Eisenstein character, here the scalar $667=1+2^{11}$ for $p=2$ in both degrees. The cone of the restriction $r$ gives the compactly
supported cohomology $H^{*}_{c}$ with its long exact sequence {\ttfamily[C,~sec:cone]}, {\ttfamily[C,~eq:coneles]}
(here $H^{0}=0$)
\[
  0\to H^{0}_{\partial}\to H^{1}_{c}\to H^{1}\xrightarrow{\,r^{1}\,}H^{1}_{\partial}\to H^{2}_{c}\to\cdots,
\]
whose exactness the snippet reads; $r^{1}$ has rank $1$ and the interior cohomology $H^{1}_{!}=\ker r^{1}$ has dimension $2$. The Hecke action on
$H^{*}_{c}$ is defined at cochain level through a \emph{pair homotopy} $\eta$ over $\Z$ {\ttfamily[C,~def:pairhomotopy]},
{\ttfamily[C,~def:coneaction]}; it is well defined and independent of $\eta$ {\ttfamily[C,~thm:conewelldef]} \PROVED. The snippet prints
$\eta$'s certificate at $p=2$ (solved over $\Z$, integral, re-derived) and the cone-action certificate.

\snippet{eisenstein_sign_cone_10_691}

\begin{lookfor}
\item \code{Eisenstein line in H\^{}1\_+: true}: the production sign; with $\iota^{*}_{\mathrm{geo}}$ it would be \code{false};
\item \code{plus : \dots loewy [1, 1] \dots cyclic} against \code{minus : \dots loewy [1]}: the extension is in the $+$ part only;
\item \code{rank r\^{}1 = 1}, \code{dim H\^{}1\_! = 2}, \code{LES exact: true}.
\end{lookfor}

%======================================================================
\part{Congruences modulo \texorpdfstring{$\ell$}{l}: local Hecke algebras}\label{part:pg-congruences}

Over $\Q$ the Hecke algebra acting on $H^{1}$ is semisimple: a product of number fields, one per Galois orbit of eigenforms (and $\Q$ for
the Eisenstein series). Modulo $\ell$ two eigensystems can coincide --- a \emph{congruence} --- and then the generalised eigenspace of the
common system $\m$ is a module over a local Artin algebra $A_\m$ which need not be a field. This part reads congruences as such local
algebras: their Hilbert functions, their Loewy series, and the extensions of Hecke modules they encode [E,~sec:ring]. It starts with the
Eisenstein--cusp congruences (Ramanujan, Kummer, Ribet), goes on to cusp--cusp congruences (visible, hidden, dihedral, multiplicity two),
the weight-raising operator $\theta$ that links congruences in different weights, and the longest local algebra found by snippet; then the
census of all local algebras up to $n=1000$ (edition~4.0).

%======================================================================
\section{Ramanujan's congruence as an Eisenstein ideal}\label{sec:pg-ramanujan}

The discriminant $\Delta=\sum\tau(n)q^{n}$ is the only normalised cusp eigenform of level one and weight $12$, and
$E_{12}=-\tfrac{B_{12}}{24}+\sum\sigma_{11}(n)q^{n}$ with $B_{12}=-\tfrac{691}{2730}$. Because $691$ divides the numerator of $B_{12}$, the
constant term of $E_{12}$ vanishes modulo $691$ and $E_{12}$ is congruent to a cusp form, necessarily $\Delta$: Ramanujan's
$\tau(n)\equiv\sigma_{11}(n)\pmod{691}$ (\LIT). This is the project's flagship example, [E,~cl:ramanujan] and [EX,~ex:691].

On cohomology, $H^{1}(\SL_2(\Z);\Sym^{10}\F_{691}^{2})$ has dimension $3$: the two classes of $\Delta$ and the Eisenstein class. All three
carry the same eigenvalue $a_p=1+p^{11}\bmod691$ of every $T_p$, so $H^{1}$ is a single generalised eigenspace, and the question is whether
the Hecke action on it is semisimple. It is not: $T_p-a_p$ is nilpotent of Jordan type $(2,1)$ for every $p$ tested, including $U_{691}$
(eigenvalue $\tau(691)\equiv1$). The congruence between the Eisenstein series and $\Delta$ is thereby a \emph{non-split extension} of Hecke
modules, and the $\F_{691}$-algebra generated by the $T_p$ on the block is $\F_{691}[\varepsilon]/(\varepsilon^{2})$: the maximal ideal
$\m=(T_2+24)=(T_2-667)$ is the \emph{Eisenstein ideal} (Mazur's term at prime level; here at level one, weight $12$), and
$\varepsilon$ measures the congruence.

The snippet derives the eigenvalues from $\sigma_{11}$ (they are output, not input), checks the Jordan types, and reads the extension
record of the cell [E,~sec:ring]: one maximal ideal, residue field $\F_{691}$, local algebra with Hilbert function $(1,1)$, Gorenstein, and the
Loewy profile of $H^{1}_\m$, $[2,1]$: a free rank-one module over $\F_{691}[\varepsilon]/\varepsilon^{2}$ plus one more line.

\snippet{ramanujan_10_691}

\begin{lookfor}
\item \code{bernoulli\_mod(12, 691) = 0}: the source of the congruence;
\item the same Jordan type $(2,1)$ for every $p$, i.e.\ one non-split $2\times2$ block and one $1\times1$ block, \emph{not} a scalar matrix;
\item \code{Hilbert function [1, 1]}: the local algebra is $\F_{691}[\varepsilon]/\varepsilon^{2}$; a reduced algebra would print \code{[1]}.
\end{lookfor}

%======================================================================
\section{Irregular primes, \texorpdfstring{$\zeta(1-k)$}{zeta(1-k)} and the Herbrand--Ribet congruence}\label{sec:pg-kummer}

The special values $\zeta(1-k)=-B_k/k$ ($k\ge2$ even) control Eisenstein congruences: if $\ell\mid B_k$ with $2\le k\le\ell-3$ (an
\emph{irregular pair} $(\ell,k)$), the constant term of $E_k$ vanishes mod~$\ell$ and $E_k$ is congruent to a cusp form of weight $k$; Ribet
used that cusp form to build an unramified $\Z/\ell$-extension of $\Q(\zeta_\ell)$ on which $\mathrm{Gal}(\Q(\zeta_\ell)/\Q)$ acts by
$\omega^{1-k}$, proving the converse of Herbrand's theorem (\LIT). Kummer's criterion says $\ell$ is irregular iff such a $k$ exists.

The console computes $B_w\bmod\ell$ for $0\le w<\ell-1$ by the $\Fl$-recursion (\code{bernoulli\_mod}, von Staudt--Clausen bounds the range).
The snippet lists the irregular indices of $37,59,67,101,103,691$ --- $32$, $44$, $58$, $68$, $24$, and $12,200$ --- and then compares the
Eisenstein block of $H^{1}(\SL_2(\Z);\Sym^{n}\F_{37}^{2})$ at the regular weight $k=12$ ($n=10$) with the irregular weight $k=32$ ($n=30$). The
Eisenstein ideal is found by its character $a_2=1+2^{k-1}$. At $k=12$ it is a one-dimensional block, nothing is glued; at $k=32$ it is
three-dimensional, with the Loewy profile $[2,1]$ and local algebra $\F_{37}[x]/x^{2}$ of \S\ref{sec:pg-ramanujan}: Ribet's cusp form of
weight $32$ is glued to $E_{32}$.

The field-theoretic side of these two cases --- the Kummer generators $X^{37}-\eta_{32}$ and $X^{691}-\eta_{12}$ over the cyclotomic
fields, the $\Z/\ell^{2}$ layers --- is in [L37,~sec:fieldsribet] and [L37,~app:ribetdata]; the $\ell$-adic $L$-values (Kubota--Leopoldt,
Coates--Wiles, the measure congruence) are [E,~sec:lvalue], [E,~sec:lvmeasure], [E,~cl:lvmeascong], and the rungs of the $\Lambda$-adic
Eisenstein ideal at $691$ are [U,~prop:u-691]. The Kummer test of the generators is the verb \code{kummer\_unit} (\S\ref{sec:pg-kummerunit}); the snippet here shows the
residue modulo $\ell$ of $\zeta(1-k)$ and its consequence on cohomology.

\snippet{kummer_ribet_30_37}

\begin{lookfor}
\item \code{ell = 691 irregular indices k: 12 200}: $691$ is irregular twice;
\item at $(10,37)$ the Eisenstein ideal has \code{dim 1}, \code{Loewy [1]}; at $(30,37)$ it has \code{dim 3}, \code{Loewy [2, 1]},
  \code{T\_m Hilbert [1, 1]};
\item the other ideal at each cell is a cusp eigenform not congruent to an Eisenstein series ($\Delta$ at $k=12$, the second weight-$32$ form at
  $k=32$), with a semisimple two-dimensional block.
\end{lookfor}

%======================================================================
\section{Kummer's criterion on cyclotomic units (\texttt{kummer\_unit})}\label{sec:pg-kummerunit}

The cyclotomic unit $\eta_i$ of the first layer of $\Q(\zeta_\ell)$ (its $\omega^{i}$-component) satisfies $v_\pi(\eta_i^{\ell-1}-1)>i$, i.e.\ it is closer
to~$1$ $\pi$-adically than its eigencharacter forces, exactly when $\ell$ divides the Bernoulli number $B_i$ (Kummer; [E,~sec:lvnf]). It is not a global
$\ell$-th power at any index shown (the record's \code{not\_lth\_power}, tested at split primes); at an irregular index the congruence is what makes the Kummer
extension $\Q(\zeta_\ell,\eta_i^{1/\ell})$ a candidate for an unramified one. The verb \code{kummer\_unit(ell, i)} returns the unit on the power basis of $\Z[\zeta_\ell]$ with the
valuation $v_\pi(\eta_i^{\ell-1}-1)$; the test is $v_\pi>i$. The record also carries the comparison with the direct Bernoulli computation
(\code{oracle\_agrees}). At $\ell=37$ the index $i=32$ is the irregular one of \S\ref{sec:pg-kummer}, at $\ell=59$ it is $i=44$ (\LIT; \CMPD).

\snippet{kummer_unit_37_59}

\begin{lookfor}
\item \code{v\_pi} equals $i$ at the regular indices and exceeds $i$ at the irregular ones ($38>32$, $60>44$);
\item \code{oracle agrees: true} on every line.
\end{lookfor}

%======================================================================
\section{A visible congruence between eigenforms}\label{sec:pg-visible}

Two eigenforms $f\ne g$ of the same weight can be congruent modulo $\ell$ (all their Hecke eigenvalues agree mod a prime above $\ell$).
On cohomology this can show up in two ways. In the notation of \S\ref{sec:pg-runs} (``Two $\ell$-adic objects''), let
$\Lambda^{\mathrm{fr}}_\m$ be the free quotient of the localisation of $H^{1}(\SL_2(\Z);\Sym^{n}\Zl^{2})$ at a residual system $\m$,
$\T_\m$ the $\Zl$-algebra generated by the $T_p$ on it, $\T_\m^{\mathrm{sat}}$ its saturation in $\mathrm{End}(\Lambda^{\mathrm{fr}}_\m)$ and
$\mathsf A_K$ the image of $\T_\m$ in $\mathrm{End}(\Lambda^{\mathrm{fr}}_\m/\ell^{K})$ {\ttfamily[C,~def:satdata]}. The mod-$\ell$ algebra
$\mathsf A_1=\T_\m/(\T_\m\cap\ell\T^{\mathrm{sat}}_\m)$ is what $H^{1}$ with $\Fl$-coefficients sees when the block has no torsion (then
$\Lambda^{\mathrm{fr}}_\m/\ell$ is the block and $\mathsf A_1=A$; on a torsion block $\mathsf A_1$ is a quotient of $A$); $\mathsf A_K/\ell$ ($K$
large) is $\T_\m/\ell$. The congruence is \emph{visible} if already $\mathsf A_1\ne\Fl$.

The console's \code{hecke\_ring\_algebra(n, ell, kmax)} runs the $\Z/\ell^{K}$ Hecke-algebra route with its certified reading
{\ttfamily[C,~cor:stablecert]}: per block it prints the Hilbert functions of $A$ (\code{A1}, on the full block) and of $\mathsf A_K/\ell$
(\code{AK/ell}, on the free part), the free rank, the torsion, and the verdict. On a torsion-free block (\code{stable}) the two are the objects
of the previous paragraph; on a block with torsion and free rank $f>0$ (\code{torsion\_free\_stable}) \code{AK/ell} is read on
$\Lambda^{\mathrm{fr}}_\m$ while \code{A1} stays the algebra $A$ of the full block, and \code{hidden} compares the two. At $(48,7)$ (weight $50$) the block with system $(a_2,a_3,a_5,a_7)=(1,0,0,0)$ has
$\mathsf A_1=\mathsf A_K/\ell=\F_7[x]/x^{2}$ (Hilbert function $(1,1)$), verdict \code{stable} (certified): two eigenforms congruent
mod~$7$, and $H^{1}$ with $\F_7$-coefficients already shows the nilpotent $x$. The second half reads the same block from the $\F_7$
extension record: a four-dimensional ideal with local algebra of Hilbert function $(1,1)$ and nilpotency index $2$. The cell is the first one
of the console's frozen script \file{029\_hecke\_ring\_algebra.oc}.

\snippet{visible_congruence_48_7}

\begin{lookfor}
\item \code{A1 [1, 1]} \emph{equal to} \code{AK/ell [1, 1]}, \code{hidden false}: visible;
\item \code{reduced\_rank1} blocks are single eigenforms (rank-one Hecke algebra); \code{torsion\_stable} is a \emph{pure-torsion} block
  ($f=0$: here $\mathbb L_\m\cong\Z/7$, the Dickson class of $n=48=6(7+1)$; $\Lambda^{\mathrm{fr}}_\m=0$ and \code{AK/ell} is read on
  $\mathbb L_\m$ itself, {\ttfamily[C,~lem:freepart]}(3)), not a block with a free part --- that verdict is \code{torsion\_free\_stable};
\item \code{nu 2} in the $\F_7$ record for the same block.
\end{lookfor}

%======================================================================
\section{A hidden congruence: \texorpdfstring{$(300,23)$}{(300,23)}}\label{sec:pg-hidden}

A congruence is \emph{hidden} if at least two eigenforms share the residual system $\m$ while $\mathsf A_1=\Fl$
([LG,~def:hiddencong], restating [EX,~s14:wp6]): $(T_p-a_p)\Lambda^{\mathrm{fr}}_\m\subseteq\ell\Lambda^{\mathrm{fr}}_\m$ for every $p$, but $T_p-a_p\notin\ell\T_\m$ for
some $p$. The congruence is then invisible on $H^{1}$ with $\Fl$-coefficients and appears only over $\Z/\ell^{K}$, $K\ge2$; its marker is
$\dim\mathsf A_1<\dim\mathsf A_K/\ell$, its depth the saturation defect $\delta=\mathrm{length}(\T_\m^{\mathrm{sat}}/\T_\m)\ge1$.

At $(n,\ell)=(300,23)$ (weight $302$) the block with system $(a_2,a_3,a_5,a_7,a_{23})=(19,14,0,0,0)$ is $\rhobar_\Delta\otimes\omega^{2}$,
dihedral for $K=\Q(\sqrt{-23})$ (\S\ref{sec:pg-dihedral}), non-ordinary ($a_{23}=0$), of dimension $D=4$: two weight-$302$ eigenforms. Mod $23$,
$\mathsf A_1=\F_{23}$ (one eigensystem, Hilbert function $(1)$); over $\Z/23^{2}$, $\mathsf A_K/\ell=\F_{23}[x]/x^{2}$ (Hilbert function $(1,1)$): the
two forms are congruent mod~$23$ but not mod $23^{2}$, depth $\delta=1$. The record's \code{hidden} flag is true, and it is the only hidden
block of the $14$ at this cell. The dihedral pattern is checked on the system itself: $a_p=0$ exactly at the primes inert in $K$. Hidden
congruences off the dihedral locus, the ``Frobenius layer'', are [E,~pr:frobeniuslayer]; the dihedral reading is [LG,~sec:cm].

\snippet{hidden_congruence_300_23}

\begin{lookfor}
\item \code{A1\_hilbert = [1]} against \code{AK\_hilbert = [1, 1]}: the congruence that $\Fl$-coefficients miss;
\item compare \S\ref{sec:pg-visible}, where both are \code{[1, 1]};
\item \code{OTHER BLOCKS hidden: 1 of 14}: hidden congruences are rare; eight other $D=4$ blocks are visible.
\end{lookfor}

%======================================================================
\section{A dihedral residual representation: Wilton's congruence mod \texorpdfstring{$23$}{23}}\label{sec:pg-dihedral}

The mod-$23$ Galois representation $\rhobar_{\Delta,23}$ is induced from $K=\Q(\sqrt{-23})$: $K$ has class number $3$, its Hilbert class field
$H$ is the splitting field of $x^{3}-x-1$, and $\rhobar_{\Delta,23}\cong\mathrm{Ind}_K^{\Q}\chi$ for a character $\chi$ of
$\mathrm{Gal}(H/K)\cong\Z/3$ (\LIT: Wilton 1930; Serre, Swinnerton-Dyer). Taking traces of Frobenius:
\[
  \tau(p)\equiv\begin{cases}0 & p\text{ inert in }K,\ (\tfrac{-23}{p})=-1,\\
  2 & p=x^{2}+xy+6y^{2}\ (\text{split, principal prime ideals}),\\
  -1 & p=2x^{2}+xy+3y^{2}\ (\text{split, non-principal}),\end{cases}\pmod{23}.
\]
The snippet reads $\tau(p)\bmod23$ off $H^{1}(\SL_2(\Z);\Sym^{10}\F_{23}^{2})$ basis-independently: since $10<23$ the reduction of the
characteristic-zero Hecke module is the $\F_{23}$-module, so $\operatorname{tr}T_p=2\tau(p)+(1+p^{11})$ (two cusp classes, one Eisenstein
class), and $\tau(p)=(\operatorname{tr}T_p-1-p^{11})/2$ in $\F_{23}$. It classifies each $p<175$, $p\ne23$, by splitting and by the principal
form (via $4p=(2x+y)^{2}+23y^{2}$) and asserts the three cases. The first principal primes are $59,101,167,173$; note that
$1+p^{11}\equiv1+(\tfrac{p}{23})$ as well, so the Eisenstein eigenvalue also vanishes at the inert primes.

Dihedral residual representations are where hidden congruences concentrate (\S\ref{sec:pg-hidden}, [LG,~sec:cm]); the dihedral ordinary
blocks of the census are [EX,~s14:dihedral].

\snippet{dihedral_wilton_10_23}

\begin{lookfor}
\item \code{tau(p) mod 23 = 0} on every inert line, \code{= 2} exactly at \code{p = 59, 101, 167, 173}, \code{= 22} ($=-1$) elsewhere;
\item the eigenvalues come from the traces of the computed $T_p$, not from a table of $\tau$.
\end{lookfor}

%======================================================================
\section{Multiplicity one fails mod \texorpdfstring{$5$}{5}}\label{sec:pg-mult2}

Over $\Q$ every Hecke eigensystem of level one occurs in $H^{1}_{\pm}$ with multiplicity one. Mod $\ell$ this fails: distinct eigenforms can share
a residual system, and they can share it \emph{including} $U_\ell$. At $\ell=5$ every residual system of level one is reducible (a $\theta$-twist of
an Eisenstein system); the atlas $\Sigma(5)$ [E,~sec:atlas] labels them by their characters $\chi=[a,b]$, $\rhobar=\omega^{a}\oplus\omega^{b}$.

At weight $36=12+5^{2}-1$ ($n=34$), $H^{1}$ has dimension $7$ ($\dim S_{36}=3$). The extension record finds three maximal ideals:
\begin{itemize}[leftmargin=1.6em]
\item atlas system~$2$, $\chi=[1,2]$, i.e.\ $a_p\equiv p(1+p)\equiv\tau(p)$ (Ramanujan's mod-$5$ system): one ideal of dimension $4$, nilpotency
  $1$ (the local algebra is $\F_5$), $U_5=0$ --- two eigenforms congruent to $\Delta$ mod~$5$ that no Hecke operator, $U_5$ included, separates,
  with the module $\F_5^{4}$ semisimple ($2$ in $H^{1}_{+}$). This is [EX,~ex:mult4]: Glover's composition factors count the multiplicity and the
  $\theta$-images from weight $30$ account for it;
\item atlas system~$1$, $\chi=[0,3]$ (that of $E_{36}$): two ideals with the same $a_2,a_3$, separated by $U_5$
  ($U_5\equiv1$ on the Eisenstein line, $U_5\equiv0$ on the two classes of a cusp form with the same anemic system).
\end{itemize}
So the anemic Hecke algebra (without $U_\ell$) does not separate the second pair, and the full Hecke algebra does not separate the first.

\snippet{multiplicity_two_34_5}

\begin{lookfor}
\item \code{ideal 2 dim 4 \dots chi [1, 2] \dots nu 1}: multiplicity two (four classes, two per sign) with a reduced algebra;
\item \code{ideal 0} and \code{ideal 1}: identical \code{a\_2}, \code{a\_3}, different \code{U\_5};
\item \code{S\_sep [2]}: $T_2$ alone separates the residual systems of $\Sigma(5)$, not the ideals.
\end{lookfor}

%======================================================================
\section{The \texorpdfstring{$\theta$}{theta}-operator and the Dickson sequence}\label{sec:pg-theta}

On mod-$\ell$ modular forms, $\theta=q\,\tfrac{d}{dq}$ raises the weight by $\ell+1$ and multiplies $a_p$ by $p$ (\LIT, Serre, Swinnerton-Dyer).
On cohomology the same map is multiplication by the Dickson invariant $L=X^{\ell}Y-XY^{\ell}$ {\ttfamily[C,~def:dicksoninv]}, which gives a short exact
sequence of $\GL_2$-modules over $\Fl$
\[
  0\to\Sym^{n}\otimes{\det}\xrightarrow{\;\cdot L\;}\Sym^{n+\ell+1}\to Q_n\to0
\]
and the map $\theta_L$ on cohomology {\ttfamily[C,~def:theta]}. Its long exact sequence {\ttfamily[C,~thm:thetales]} \PROVED{}
\[
  0\to H^{0}_{n}\xrightarrow{\theta}H^{0}_{\uparrow}\to H^{0}(Q_n)\xrightarrow{\delta}H^{1}_{n}\xrightarrow{\theta}H^{1}_{\uparrow}\to H^{1}(Q_n)\to0
\]
is exact, so $\ker\theta^{1}=\operatorname{im}\delta$; and $\theta$ is Hecke-equivariant only up to the determinant twist,
$T_p\,\theta=p\,\theta\,T_p$, with $U_\ell\,\theta=0$ {\ttfamily[C,~prop:thetatwist]} \PROVED. The console objects are
[E,~sec:thetaobj].

At $(16,7)$, $n+\ell+1=24$ (weights $18\to26$): the sequence is exact with $\delta$ of rank $1$ (a Dickson class of $H^{0}(Q_{16})$ kills one class of
$H^{1}$), $\theta^{1}$ has rank $3$, and the twisted identity holds at $p=2,3,5$ while the untwisted one fails. The eigenvalues of $T_p$ on the two
$H^{1}$'s show the twist directly: $\{4,5\}\cdot2=\{1,3\}$, $\{0,1,6\}\cdot3=\{0,3,4\}$, $\{0,3,4\}\cdot5=\{0,1,6\}$ in $\F_7$ (the images
match because $\theta^{1}$ has rank $3$ and both spaces have dimension $4$).

\snippet{dickson_theta_16_7}

\begin{lookfor}
\item \code{ranks \{theta0: 1, pi0: 0, delta: 1, theta1: 3, pi1: 1\}} and \code{exact: true};
\item \code{A\_up theta == p theta A: true} against \code{A\_up theta == theta A: false};
\item multiply the left eigenvalue set by $p$ modulo $7$ to get the right one.
\end{lookfor}

%======================================================================
\section{A mod-\texorpdfstring{$37$}{37} Hecke algebra of Loewy length \texorpdfstring{$4$}{4}}\label{sec:pg-loewy4}

The census found the first mod-$\ell$ Hecke algebras of Loewy length $4$ on $H^{1}$ at $\ell=37$ (edition~4.0 adds $\ell=67$ and $\ell=139$,
\S\ref{sec:pg-census34}), at the residual representations
$\rhobar_a=\omega^{a}\oplus\omega^{a+5}$ from weight $k=222+38a$ on [L37]: the local algebra is $\F_{37}[x]/x^{4}$ and the $8$-dimensional block is
$A\oplus A/x^{3}\oplus A/x$. The block is cut out by four cusp eigenforms reducing to $\theta^{a}E_6\equiv\theta^{a+5}E_{32}$, and [L37]'s dimension formula traces the rank
$4$ to Ribet's cusp form of weight $32$ (\S\ref{sec:pg-kummer}); the families and the fields are [L37,~sec:families],
[L37,~sec:fieldsribet].

The snippet takes the first cell of the $a=1$ family, $(258,37)$, weight $260$. Part~(1) is the $\F_{37}$ extension record (two seconds): the
$8$-dimensional ideal with $a_2=2+2^{6}=29$, $U_{37}=0$ (non-ordinary), local algebra of Hilbert function $(1,1,1,1)$, Gorenstein, embedding
dimension $1$, Loewy series $[3,2,2,1]$ and Jordan type $(4,3,1)$ of the nilpotent generator: $H^{1}_\m\cong A\oplus A/x^{3}\oplus A/x$. Part~(2) is the
$\Z/37^{2}$ route (about a minute): the block is torsion-free of rank $8$, $\mathsf A_1=\mathsf A_K/\ell$ with Hilbert function $(1,1,1,1)$, so
nothing is hidden and the census algebra is the saturated one, by the certified reading {\ttfamily[C,~cor:stablecert]}.

\snippet{loewy4_258_37}

\begin{lookfor}
\item \code{Hilbert function [1, 1, 1, 1]}: $\F_{37}[x]/x^{4}$, Loewy length $4$;
\item \code{Jordan type of n\_2 [4, 3, 1]}: the module $A\oplus A/x^{3}\oplus A/x$, not cyclic;
\item \code{A1 [1, 1, 1, 1]} $=$ \code{AK/ell [1, 1, 1, 1]}, \code{verdict stable (certified)}.
\end{lookfor}

%----------------------------------------------------------------------
\section{Census 4.0 at a glance: the shape of every non-reduced algebra (runs)}\label{sec:pg-census34}\label{sec:pg-census40}

\paragraph{Six algebras.} \CMPD{} (each row certified) and \EMP{} (as a pattern). Every one of the $49\,176$ non-reduced census algebras of
the signed edition~4.0 (main $+$ b1) is one of six algebras ($\F_q$ the residue field):
\begin{center}\small
\begin{tabular}{@{}lrll@{}}
\toprule
$A$ & rows & Hilbert function & Gorenstein \\
\midrule
% PG-CENSUS-SHAPES-BEGIN (read by docs/latex/tools/pg_census_check.py: one row per shape, `$A$ & rows & $(h)$ & yes|no`)
$\F_q[x]/x^{2}$ & $46\,185$ & $(1,1)$ & yes \\
$\F_q[x]/x^{3}$ & $414$ & $(1,1,1)$ & yes \\
$\F_q[x]/x^{4}$ & $248$ & $(1,1,1,1)$ & yes \\
$\F_q[x,y]/(x,y)^{2}$ & $2016$ & $(1,2)$ & no \\
$\F_q[x,y,z]/(x,y,z)^{2}$ & $303$ & $(1,3)$ & no \\
$\F_q[x,y]/(x^{3},xy,y^{2})$ (embedding dimension $2$, $\m_A^{3}=0\ne\m_A^{2}$) & $10$ & $(1,2,1)$ & no \\
% PG-CENSUS-SHAPES-END
\bottomrule
\end{tabular}
\end{center}
The shape $(1,2,1)$ is a single algebra: $\m_A^{2}$ is a line, the multiplication form on $\m_A/\m_A^{2}$ is non-degenerate exactly when $A$ is
Gorenstein, and a non-Gorenstein one has a rank-one form, unique up to scaling, whence $\F_q[x,y]/(x^{3},xy,y^{2})$.
Section b1 contributes only $(1,1)$ ($2542$ rows) and $(1,1,1,1)$ ($15$ rows); the partial section b2 (pre-release) has $825$ non-reduced
rows, all $\F_q[x]/x^{2}$. So, on the signed census:
\begin{itemize}[leftmargin=1.6em,itemsep=1pt]
\item a census algebra is Gorenstein if and only if it is curvilinear ($0$ exceptions);
\item the embedding dimension is at most $3$, the tangent-space bound of [LG,~prop:tangent];
\item a non-curvilinear algebra has $\m_A^{2}=0$, and an embedding dimension of at least $2$ never comes with Loewy length $\ge3$ --- \emph{except} on
  the $10$ rows of shape $(1,2,1)$, all at $n>700$: eight dihedral rows at $\ell=23$ ($n=826,\dots,980$, $D=16$, module Loewy series $[8,4,4]$;
  \S\ref{sec:pg-dihedral3}) and two at $\ell=59$, $(942,59)$ and $(1000,59)$ ($D=8$, $[4,2,2]$). On edition~3.4's window $n\le700$ both
  statements hold without exception.
\end{itemize}

\paragraph{Where each shape occurs.}
\begin{itemize}[leftmargin=1.6em,itemsep=1pt]
\item \textbf{Loewy length $\ge3$} occurs on $672$ rows, at $\ell\in\{23,31,37,59,67,139,547\}$:
  \begin{itemize}[itemsep=0pt]
  \item $418$ dihedral rows at $23$ and $31$, Loewy length $3$ (\S\ref{sec:pg-dihedral3});
  \item $233$ Kummer--Frobenius rows of Loewy length $4$: $222$ at $37$ from $n=258$ (\S\ref{sec:pg-loewy4}, [L37]) and $11$ at $67$ from
    $n=736$, at the residual representations $\omega^{a}\oplus\omega^{a+9}$ ($67$ is irregular at $k=58$, and $9=67-58$ as $5=37-32$); all
    listed in \pfile{kummer_frobenius40.tsv};
  \item two Kummer rows at $547$: $(484,547)$, which is ordinary (\S\ref{sec:pg-ord547}), and $(608,547)$;
  \item $4$ rows at $59$ from $n=884$, Loewy length $3$, labelled $\theta^{0}(f_{16})$, $\theta^{1}(f_{16})$ (the label rests on evidence at
    $61$ primes, not on a certificate);
  \item $15$ rows at $139$ (section b1) from $n=416$: $A=\F_{139}[x]/x^{4}$, module Loewy series $[2,2,2,2]$, $D=8$, labelled
    $\theta^{j}(f_{138})$, $j=2,\dots,6$; they are not Kummer--Frobenius.
  \end{itemize}
  The rows at $59$ and $139$ are new in 4.0 (at $n>700$, resp.\ at the new primes). A referee check against the census records
  (2026-11-06; four-prime comparison of the systems $(a_2,a_3,a_5,a_7)$, not a certificate) gives: at $\ell=139$ the class has a congruence in
  weight $138$ (the ideal $(136,139)$ id~$2$, $D=4$, Hilbert function $(1,1)$; evidence on four primes), so it is not congruence-free and
  [E,~pr:latticesplit] (R) says nothing about these $15$ rows. At $\ell=59$ the class is irreducible, of residue degree one, not dihedral
  ($(2/59)=-1$ and $a_2\ne0$) and congruence-free in the sense of \S\ref{sec:pg-dihedral3} (the only low-weight matches are single forms, at
  weights $16$ and $46$), so the $4$ rows at $59$ are a \emph{candidate} counterexample to (R), under independent verification and pending the
  owner's ruling; the gallery claims no more than this. (1.4.2: verified independently, \pfile{docs/implementation/slot25/R59_verification.md},
  $(T_2-a_2)^{2}\ne0$ recomputed at all four cells; the class is non-ordinary, $a_{59}\equiv0$. By owner ruling of 2026-11-06 [E] 1.25.0-pre.1
  restates (R): on irreducible classes it assumes ``of one weight'', exactly one pair $(w,j)$ of the test with a non-zero eigenspace, which the
  class at $59$, at $(16,0)$ and $(46,14)$, is not.)
\item \textbf{The non-Gorenstein algebras} occur only at $\ell\in\{7,11,13,17,19,23,29,59\}$. Their first cells are $n=82$, $174$, $206$,
  $338$, $416$, $596$, $926$ at $\ell=7$, $11$, $13$, $17$, $19$, $23$, $29$ ($421$, $680$, $569$, $282$, $222$, $148$, $5$ rows), roughly
  $n\approx\ell^{2}+30\dots85$ (\EMP); at $59$ they are the two $(1,2,1)$ rows above, from $n=942$, off that pattern.
\item \textbf{Embedding dimension $3$} occurs only at $\ell=7$, $11$, $13$, $17$, from $n=142$, $354$, $496$, $852$.
\end{itemize}
In [EX] these first appear at $(82,7)$ [EX,~ex:nongor] and $(142,7)$ [EX,~ex:largest]. That document's ``What we have not seen'' list
[EX,~sec:notseen] is scoped to $n\le200$ and remains true there. ``Only at $\ell\le29$'' (non-Gorenstein, apart from the two rows at $59$) and
the list of primes with Loewy length $\ge3$ are statements about the window of the census (even $n\le1000$, $\ell\le199$): the onsets grow like
$\ell^{2}$ (non-Gorenstein, dihedral) or $\ell w$ (Kummer--Frobenius), and the window of~4.0 already moved both lists.

\paragraph{The literature} (search of 2026-10-19). Failure of the Gorenstein property for mod-$p$ Hecke algebras is classical in weight~$2$
and weight~$p$, through ``unramified at $p$ with scalar Frobenius'' (Kilford--Wiese \cite{ph_kilford_wiese}); the embedding-dimension bound is
the tangent-space bound (Bella\"iche--Khare \cite{bk}, Deo \cite{deo}). Our $A$ is the image on mod-$\ell$ cohomology; when that agrees
with $\T_\m/\ell$ is the subject of Wiese \cite{ph_wiese_faithful}. For the dichotomy ``curvilinear or $\m^{2}=0$'' on edition~3.4 (which
4.0 breaks at $n>700$ by the ten $(1,2,1)$ rows) and for the non-Gorenstein onset near $n\approx\ell^{2}$ in weights $k\gg\ell+1$, \emph{no
prior was found in our search (2026-10-19)}. A mechanism to test: the first $n$ at which a Serre weight occurs twice among the composition
factors of $\Sym^{n}\Fl^{2}$.

\paragraph{Hidden congruences} ($\dim A<\dim\T_\m/\ell$ on a torsion-free block; \S\ref{sec:pg-hidden}) are certified on $196$, $338$, $164$,
$144$, $31$, $10$ rows at $\ell=7$, $11$, $13$, $17$, $19$, $23$ ($883$ rows; the flag is computed only where a weight-shift run exists, on
$16\,117$ rows, and never on a block with torsion; a further $459$ rows with a run read \code{not\_stable}: $454$ torsion blocks and $5$ torsion-free blocks not stable by \code{KTOP}). None occurs on the computed rows at any $\ell\ge29$, section b1 included ($738$ rows read
\code{false}). At $17$ and $19$ they begin in ``Frobenius layers'' $n=\ell k_0+j(\ell+1)-2$ [E,~pr:frobeniuslayer]. At $23$ all ten are
dihedral twists of $\Delta$.

\noindent Sources: \pfile{docs/matrix-engine/results/congruence_census/edition3/4.0/} (\pfile{rows40.tsv.gz}, \pfile{rows40_b1.tsv.gz},
\pfile{rows40_b2.tsv.gz}, \pfile{category_table40*.txt}, \pfile{findings40*.txt}, \pfile{kummer_frobenius40.tsv}); every number recomputed by
\file{docs/latex/tools/pg\_census\_check.py}; the onset analysis \pfile{docs/implementation/slot22/C33_followups_report.md} \S F33-3.

%----------------------------------------------------------------------
\section{Dihedral classes: Loewy length 3 on congruence-free classes (runs)}\label{sec:pg-dihedral3}

\paragraph{The conjecture.} The conjecture [E,~pr:latticesplit] (``nearly split lattice'', (R)) predicts $\m_A^{2}=0$ on every congruence-free
class: a class whose Serre weights $w\le\ell+1$ carry no congruence (every generalised eigenspace in $M_w(\Fl)$ has dimension $\le1$).

\paragraph{The counterexample.} Census~3.3 contradicted it on $161$ rows (\CMPD; report F33-2, refereed, owner ruling of 2026-10-09):
\begin{itemize}[leftmargin=1.6em,itemsep=1pt]
\item \emph{The classes.} They are the twists of $\rhobar_\Delta$ mod $23$ and of $\rhobar_{f_{16}}$ mod $31$. Both are dihedral, induced
  from $\Q(\sqrt{-\ell})$, and congruence-free. Their eigensystems are $(a_2,a_3,a_5,a_7)\equiv(22,22,0,0)\bmod23$ and
  $(30,0,30,30)\bmod31$.
\item \emph{The algebra.} On those $161$ rows of 3.3/3.4 ($n\le700$) the census algebra is $A\cong\Fl[x]/x^{3}$, with module Loewy series
  $[4,2,2]$ or $[6,4,2]$: $140$ rows at $\ell=23$ (from $n=298$; $D=8$ or $12$) and $21$ at $\ell=31$ (from $n=526$). The signed census~4.0
  (main section, $n\le1000$) has $418$ dihedral rows of Loewy length $3$: $286$ at $\ell=23$ and $132$ at $\ell=31$ (onsets unchanged). Above
  $n=700$ the blocks at $23$ grow: $D=12$ with module Loewy series $[8,2,2]$ (from $n=782$) and $D=16$ with $[8,4,4]$ (from $n=826$), and on
  eight of the $D=16$ rows ($n=826,\dots,980$) $A$ is no longer curvilinear but has Hilbert function $(1,2,1)$ (\S\ref{sec:pg-census34}); every
  other dihedral row has $A\cong\Fl[x]/x^{3}$.
\item \emph{The mechanism.} $U_\ell=0$ on these blocks. $(T_q-a_q)(T_{q'}-a_{q'})\ne0$ on $H^{1}_\m$ exactly when $q$ and $q'$ both split in
  $\Q(\sqrt{-\ell})$. The inert primes and $U_\ell$ lie in $\m_A^{2}$.
\item \emph{The forms side.} $\T^{\mathrm f}_\m/\ell\cong\Fl[x]/x^{4}$ at $(298,23)$, $(540,23)$ and $(526,31)$, with $\delta=1,2,1$.
\item \emph{The $\ell$-adic signature} (mod $\ell^{6}$). The Newton slope is $1/2$ at split primes and $1$ at inert primes; there is no
  $\Zl$-rational eigenform.
\end{itemize}

\paragraph{After the counterexample.} [E] now assumes ``congruence-free \emph{and not dihedral}''. On the remaining $100\,025$
residue-degree-one, non-Dickson rows of~3.3 there is no row of Loewy length $3$. The $84$ rows of Loewy length $4$ of~3.3 are all
Kummer--Frobenius, hence never congruence-free. Census~4.0 adds Loewy-length-$\ge3$ rows that are neither dihedral nor Kummer--Frobenius:
$4$ at $\ell=59$ ($n\ge884$) and $15$ at $\ell=139$ (section b1), \S\ref{sec:pg-census34}. A referee check on the census records (\S\ref{sec:pg-census34}; four-prime evidence, not a certificate) finds that
the class at $\ell=139$ has a congruence in weight $138$, so (R) says nothing there, while the class at $\ell=59$ is congruence-free and not
dihedral: its $4$ rows are a \emph{candidate} counterexample to the restated hypothesis, under independent verification and pending the owner's
ruling. Until that is settled, the restated (R) is tested with certainty only on~3.3. (1.4.2: settled. The four rows are a counterexample, verified
from scratch, \pfile{docs/implementation/slot25/R59_verification.md}; the class is non-ordinary and has two cuspidal Serre weights. By owner ruling of
2026-11-06 [E] 1.25.0-pre.1 replaces ``not dihedral'' by ``of one weight'' on irreducible classes; that excludes the dihedral, the non-ordinary and
the companion classes, and no row of census~4.0 contradicts the new form, [E,~pr:latticesplit].)

\paragraph{The onset.} The onsets $n_0=298,526$ fit $n_0=\ell(\ell+3)/2-1$ (\EMP, two points). The next primes $\ell\equiv3\pmod4$ with
$h(-\ell)>1$ are $47$ ($h=5$), $59$, $71$, $79$, $83$, $103$, $107$, so the fit predicts $n_0=1174$ at $\ell=47$ and $1828$ at $\ell=59$, both
beyond the census grid.

\paragraph{The literature} (search of 2026-10-19): \emph{no prior found}. The nearest results are Deo's non-reduced dihedral algebras
\cite[Prop.~27--28]{deo}, whose nilpotents come from diamond operators with $p$ unramified in the CM field (here $\ell$ is ramified in
$\Q(\sqrt{-\ell})$ and there are no diamond operators), and Bella\"iche--Khare's $\F[[x,y]]$ for the algebra of all weights \cite{bk}. A
fixed-weight quotient of $\F[[x,y]]$ can be non-reduced without any obstruction, so the phenomenon may still follow from \cite{bk,deo}
and the CM geometry; the next check is $\dim H^{1}(G_{\Q,\{\ell\}},\mathrm{ad}^{0}\rhobar)$ for $\rhobar_\Delta$ mod~$23$.

\noindent Sources: \pfile{docs/implementation/slot22/C33_followups_report.md} \S F33-2, \pfile{F33-2_referee.md}, \pfile{owner_ruling_F33-2_1009.md}.

%----------------------------------------------------------------------
\section{The saturation defect in the census (runs)}\label{sec:pg-satdef}

$\delta>0$ means that the Hecke algebra $\T_\m$ on the free lattice $\Lambda^{\mathrm{fr}}_\m$ is not saturated: $\mathsf A_1=\T_\m/\ell$ fails,
and some congruence is hidden or deepened (\S\ref{sec:pg-hidden}). The census computes $\delta$ (column \code{sat\_length}), together with
$\dim\mathrm{Tor}_1^{\Zl}(\T^{\mathrm{sat}}_\m/\T_\m,\Fl)$ and its Loewy length, by \code{HeckeSaturationDefectValidation} on the torsion-free
blocks of the cells $n\le700$ of the main section; each value is certified final (a plateau of $\varepsilon_K$, {\ttfamily[C,~prop:satdefect]}(2)).
The $1593$ rows of blocks with torsion are refused by name (\S\ref{sec:pg-runs}, ``Two $\ell$-adic objects''); the cells $n>700$ and section b1
had no saturation pass. In the signed census~4.0 (\CMPD):
\begin{itemize}[leftmargin=1.6em,itemsep=1pt]
\item \emph{Coverage.} $\delta$ is computed on $62\,021$ rows (edition~3.4: $56\,287$; the step SATC of 3.5 added the cells at $\ell=5$, $7$,
  $11$, $13$, $547$, $691$, at $\ell=7$ up to $K=16$).
\item \emph{Where $\delta>0$.} $8954$ rows have $\delta>0$, all at $\ell\le53$ except one: $(272,547)$, ideal~$0$, $\delta=1$. Only $\ell\le23$
  has $\delta\ge2$.
\item \emph{The largest value.} The largest is $\delta=190$ at $(698,7)$ ideal~4, where $\mathrm{Tor}_1^{\Zl}(\T^{\mathrm{sat}}/\T_\m,\Fl)$
  has dimension $26$ and Loewy length $6$. (Edition~3.4, whose $\ell=7$ values stopped at $n=300$, had its maximum $\delta=29$ at $(298,7)$
  ideal~5, Tor of dimension $9$ and Loewy length $4$.)
\item \emph{The block of $\Delta$ at $(700,7)$} (ideal~5, $D=58$). Over $\Z/7^{K}$ the defect is $\varepsilon_K=26,49,70,91,109,125$ for
  $K=1,\dots,6$, still rising at $K=6$ (the run of slot~23); the SATC run went to $K=16$ and found the plateau: $\delta=182$, with Tor of
  dimension $26$ and Loewy length $6$.
\item \emph{The second-workstation run.} On the new cells of~3.3 ($\ell=17,\dots,59$, $n=122,\dots,700$; rows of source \code{c1a\_local}), $5060$ rows in
  $723$ cells are non-zero. The values (dimension, Loewy length, $\delta$) reach $(4,3,6)$. The non-zero cells are $219$, $195$, $158$, $75$,
  $55$, $15$, $4$, $2$ at $\ell=17$, $19$, $23$, $29$, $31$, $37$, $41$, $43$, with none at $\ell\ge47$. These rows are unchanged in~4.0.
\end{itemize}
On the forms side, [LG,~sec:where] reports $\delta=129$ for $\T^{\mathrm f}_\m$ at $(604,7)$; the census value of the same block (ideal~5,
$D=50$) on the cohomological lattice is $\delta=128$ --- two different algebras ({\ttfamily[C,~prop:cohforms]}), so the two numbers need not
agree. That the defect shrinks with $\ell$ is \EMP{} and a statement about the window $n\le700$. Context from the literature: Jochnowitz
\cite{jochnowitz} and Calegari--Emerton \cite[Thm.~2.1]{ph_ce_largeindex} (the index of $\T$ in its normalisation grows at least linearly in
$k$). Our saturation lies between $\T$ and the normalisation, so their bound makes large $\delta$ unsurprising without bounding it below; for the
census distribution itself no prior was found in our search (2026-10-19).

\noindent Sources: the census files (\pfile{rows40.tsv.gz}, columns \code{sat\_defect\_dim}, \code{sat\_defect\_loewy}, \code{sat\_length},
\code{sat\_kmax}; \pfile{../3.5/sat_kmax35_cells.tsv}), recomputed by \file{pg\_census\_check.py};
the merge report of the second workstation's saturation run (\S5), \pfile{slot22/census34_sat_ingest_report.md},
\pfile{slot23/KIT_redraft_report.md} \S5.1 (the $\varepsilon_K$ of $(700,7)$ to $K=6$).

%======================================================================
\part{Level \texorpdfstring{$\Ga_0(N)$}{Gamma0(N)} and characters}\label{part:pg-level}

Everything so far was level one. By Shapiro's lemma the cohomology of the congruence subgroup $\Ga_0(N)$ is that of $\SL_2(\Z)$ with
coefficients induced from $\Ga_0(N)$, so the same seven-vertex model computes it {\ttfamily[C,~sec:levelN]}; characters $\chi$ of
$(\Z/N)^{\times}$ twist the coefficients {\ttfamily[C,~sec:nebentypus]}. New phenomena: old forms coming from lower level, a parity
obstruction for characters, and Mazur's Eisenstein ideal of prime level --- here at the primes $2$ and $3$, where it is largest. (Its
integral face, the congruence prime $5$ at level $11$, is in \S\ref{sec:pg-descent-congruence}.)

%======================================================================
\section{Level $\Ga_0(N)$: old forms at level $2$ (\texttt{level\_n})}\label{sec:pg-leveln}

\code{level\_n(N, n, ell)} computes $H^{1}(\Ga_0(N);\Sym^{n}\Fl^{2})$ by Shapiro's lemma on the fixed model {\ttfamily[C,~sec:levelN]},
with the Heilbronn operators $T_p$ ($p\nmid N$) and $U_p$ ($p\mid N$). For $N=2$ the group has index $3$, and the weight-$12$ cohomology at
$\ell=7$ has dimension $6=2\cdot3$: $\Delta$ contributes through its two old forms $f(z)$ and $f(2z)$ (two classes each), and the Eisenstein part is that
of the two cusps (\CMPD). On the old part the operators $T_3$ and $T_5$ of level $2$ act as those of level $1$.

\snippet{level_n_old_forms_2_7}

\begin{lookfor}
\item \code{h1 = 3} at level $1$ and \code{h1 = 6} at level $2$, \code{index 3};
\item the printed \code{charpoly} lists are the coefficients from low to high degree; the factor $x^{4}$ of $T_3$ and $T_5$ at level $2$ is the four old
  classes of $\Delta$, on which they act as $\tau(3)\equiv\tau(5)\equiv0\pmod 7$, and the rest is the Eisenstein part.
\end{lookfor}

%======================================================================
\section{Characters and the parity obstruction (\texttt{nebentypus})}\label{sec:pg-nebentypus}

\code{nebentypus(N, m, n, ell)} replaces $\Sym^{n}$ by $\Sym^{n}\otimes\chi$ for the Dirichlet character of Conrey label $N.m$ reduced into $\Fl$
({\ttfamily[C,~sec:nebentypus]}); $\ell\nmid N$ and $\mathrm{ord}\,\chi\mid\ell-1$. Three checks at $N=7$, $\ell=13$, $n=2$ (\CMPD): the trivial character $7.1$
reproduces \code{level\_n} operator by operator; the cubic character $7.2$ is even, so $\Sym^{2}$ is allowed, and gives $h^{1}=4$; the sextic character $7.3$ is odd and the
cohomology vanishes ($\chi(-1)=(-1)^{n}$ is required).

\snippet{nebentypus_7_13}

\begin{lookfor}
\item \code{parity\_ok: false} and \code{h1 = 0} for $7.3$: the parity obstruction, not a computational accident;
\item the characteristic polynomial of $T_2$ for $7.2$ has degree $4$ and differs from the trivial-character one.
\end{lookfor}

%----------------------------------------------------------------------
\section{Mazur's Eisenstein ideal at \texorpdfstring{$p=2,3$}{p=2,3}, and a class-group correction}\label{sec:pg-mazur23}

\paragraph{The setting.} For a prime $N$ and $p$ prime, let $W$ be the Eisenstein-localised Hecke module on
$H_1(X_0(N),\text{cusps};\F_p)$ in weight $2$. Its Loewy length measures Mazur's Eisenstein ideal mod~$p$.

\paragraph{The construction at $p=2,3$.} The model's construction of $W$ through Manin symbols is valid at $p=2,3$ as well
{\ttfamily[C,~prop:levelNcurve]} (\PROVED).
\begin{itemize}[leftmargin=1.6em,itemsep=1pt]
\item \emph{The largest case.} At $(N,p)=(21\,929,2)$: $W$ has dimension $255$ and Loewy length $128$, with layers $255,253,\dots,1$
  (curvilinear), and $h(-4\cdot21\,929)=256$ (\CMPD). This is Calegari--Emerton's own example: $e_2=2^{m-1}-1$ for
  $N\equiv9\pmod{16}$ \cite[Thm.~1.1]{calegari_emerton}, and \S6 there states $\mathrm{Cl}(\Q(\sqrt{-21929}))\cong\Z/256$ and $e_2=127$. Our
  computation is an independent check on cohomology, with the layer structure.
\item \emph{Large $e_3$.} For $N\le10^{4}$ the cells with $e_3\ge3$ (the Eisenstein-ideal exponent at $p=3$) are:
  \begin{itemize}[itemsep=0pt]
  \item $e_3=3$ at $1747$, $2143$, $2467$, $3583$, $4177$, $4519$, $4789$, $5437$, $5527$, $8011$, $8389$, $8677$;
  \item $e_3=4$ at $6733$, $7219$;
  \item $e_3=5$ at $3061$.
  \end{itemize}
  Calegari--Emerton mention only $2143$; for this list no prior was found in our search (2026-10-19).
\item \emph{Generation without a good prime.} Mazur's good-prime criterion is sufficient, not necessary, at $p=2$: at $(2473,2)$ the
  operator $\eta_2$ generates although no prime $\le13$ is good {\ttfamily[C,~rmk:levelNcurvefamily]}.
\end{itemize}

\paragraph{A correction to Calegari--Emerton} (\CMPD, exact $h^{-}$ via $B_{1,\chi}$). Let $K^{(0)}=\Q(\sqrt{-3})K_3$ be the cyclic sextic
CM field of conductor $3N$. Calegari--Emerton \cite[\S6]{calegari_emerton} name $N=2143$ as the smallest prime $N\equiv1\pmod9$ for which $\mathrm{Cl}(K^{(0)})$ has
an element of order $9$ (their sentence was re-read verbatim, and $1297$ re-checked independently, on 2026-10-19; no published correction was
found). That is not so:
\begin{itemize}[leftmargin=1.6em,itemsep=1pt]
\item $N=1297$ already gives $\mathrm{Cl}(K^{(0)})[3^{\infty}]\cong\Z/9\oplus\Z/9$;
\item within $N\equiv10,19\pmod{27}$ the first is $1747$ ($\Z/9\oplus\Z/3$).
\end{itemize}
A charitable reading reconciles the sentence: $2143$ is the smallest $N\equiv1\pmod9$, $N\not\equiv1\pmod{27}$, with
$\mathrm{Cl}(K^{(0)})[3^{\infty}]\cong\Z/9\oplus\Z/9$.
The $3$-part has the structure $\Z_3[\omega]/\lambda^{m}$ with $m=v_3(h^{-})$ (\PROVED). Two equivalences are observed (\EMP):
\begin{itemize}[leftmargin=1.6em,itemsep=1pt]
\item $m\ge2\iff\bigl(\tfrac{N-1}{3}\bigr)!$ is a cube mod~$N$, on all $203$ primes $N\le10^{4}$, $N\equiv1\pmod9$. The criterion is
  Calegari--Emerton's \cite[Thm.~1.3(i)]{calegari_emerton} for the $3$-part of $\mathrm{Cl}(\Q(\sqrt{-3},N^{1/3}))$ (with Elkies's formula
  $((N-1)/3)!^{3}\equiv\pi\bmod\bar\pi$; see also Gerth \cite{ph_gerth05}); its form for $K^{(0)}$ was not found stated and is probably a
  reflection argument;
\item $e_3\ge2\iff m\ge2$, on the $39$ of them where $e_3$ is computed ($N\le2500$).
\end{itemize}

\paragraph{Class groups of $\Q(N^{1/5})$.} Calegari--Emerton's Lemma~5.13 holds on all $12$ cells with $g_5\ge2$ ($N\le800$). The $5$-rank is
$3$ at $N=211$, $281$, $691$. The theory is known (Chevalley, Gerth, Calegari--Emerton); these cells are still to be checked against the
$5$-rank criterion of Schaefer--Stubley \cite{ph_schaefer_stubley}.

\noindent Sources: \pfile{docs/implementation/slot22/LN23_report.md}, \pfile{LN23F_report.md}, \pfile{NFQ1_report.md},
\pfile{NF-PRANK_report.md}; \pfile{docs/phenomenology/literature_search_B.md} \S4, \S6.

%======================================================================
\part{Hida theory: ordinary families across weights}\label{part:pg-hida}

So far the weight was fixed. Hida theory moves it: the \emph{ordinary} part of $H^{1}$ --- where $U_\ell$ acts invertibly --- has the same
rank at all weights in a class modulo $\ell-1$, and its Hecke eigenvalues vary $\ell$-adically continuously, so that weights congruent modulo
$(\ell-1)\ell^{m-1}$ give characteristic polynomials congruent modulo $\ell^{m}$ {\ttfamily[C,~thm:hidacontrol]}, {\ttfamily[C,~thm:hidaweight]}.
The first section shows this on cohomology modulo $13$ and $13^{2}$; the next three show the same computation through the dedicated verbs,
on cohomology and on forms; the run sections follow one family to $n=6910$ and record the exceptional ordinary block at $\ell=547$.

%======================================================================
\section{Hida's control theorem, mod \texorpdfstring{$13$}{13} and mod \texorpdfstring{$13^{2}$}{13 squared}}\label{sec:pg-hida}

For an endomorphism $U$ of a finite $\Z/\ell^{e}$-module $M$ the \emph{ordinary projector} is $e_{\mathrm{ord}}=\lim_j U^{j!}$ and
$e_{\mathrm{ord}}M$ the ordinary part {\ttfamily[C,~def:ordinary]}, {\ttfamily[C,~prop:ordinary]}; for $H^{1}$ and $U=U_\ell$ it is the
ordinary part in Hida's sense. Hida's theorems {\ttfamily[C,~thm:hidacontrol]}, {\ttfamily[C,~thm:hidaweight]} \PROVED{} (from Hida's papers,
locators verified there) give: for $\ell\ge5$ the $\Zl$-rank of the ordinary part depends only on the weight modulo $\ell-1$, so
$h^{1}_{\mathrm{ord}}(n)=2\,\mathrm{rank}\,S^{\mathrm{ord}}_{n+2}+1$ is constant along a weight class; and the ordinary Hecke eigenvalues are
$\ell$-adically analytic in the weight, so two weights $k\equiv k'\pmod{(\ell-1)\ell^{m-1}}$ have ordinary characteristic polynomials congruent
modulo $\ell^{m}$. The project's data show the congruence level --- the minimum over the operators $T_p$, $p\le23$, and $U_\ell$ --- is \emph{exactly}
$m=1+v_\ell(\Delta k/(\ell-1))$ on every pair tested [EX,~s14:wshida]. (A single $T_p$ can agree further: on the Eisenstein factor
$1+p^{k-1}$ whenever $p^{\ell-1}\equiv1\pmod{\ell^{2}}$, e.g.\ $p=3$, $\ell=11$.) the $\Lambda$-adic Hecke algebra itself is [U,~sec:u-lambda].

$\Delta$ is $13$-ordinary ($\tau(13)=-577738\not\equiv0$). Part (1) of the snippet runs along $n\equiv10\pmod{12}$, $n=10,\dots,166$: $\dim H^{1}$
grows from $3$ to $29$, the ordinary rank (the rank of $U_{13}^{\,d}$, $d=\dim H^{1}$) stays $3$ (Eisenstein and the two classes of the $\Delta$-family),
and the characteristic polynomial of $T_2$ on the ordinary part is the same at every weight. Part (2) switches the coefficient ring to $\Z/13^{2}$
(\code{set coefficient\_ring = "Z/ell\^{}2"}, a mathematical setting, [E,~sec:ring]); $H^{1}$ is free at the three cells, the ordinary idempotent is
computed as $E=(U^{13^{8}})^{13^{12}-1}$ and certified by $E^{2}=E$ (on the non-ordinary part $U$ is topologically nilpotent; on the ordinary part
$U^{13^{8}}$ is the Teichm\"uller part of a unit, of order dividing $13^{f}-1$), and the snippet prints
$\mathrm{charpoly}(T_2E)=X^{d-3}\,f_{\mathrm{ord}}(X)$. Weights $12$ and $168$ differ by $156=12\cdot13$, so $m=2$: the ordinary factors agree
modulo $169$. Weights $12$ and $24$ differ by $12$, so $m=1$: the factors agree modulo $13$ ($27\equiv92$, $75\equiv153$, $72\equiv163$) but not modulo
$169$. \CMPD.

\snippet{hida_control_13}

\begin{lookfor}
\item \code{ordinary rank = 3} on every line of part (1), against a growing \code{dim H\^{}1};
\item the ordinary factor $X^{3}+27X^{2}+75X+72$ at \code{n = 10} and \code{n = 166}, but $X^{3}+92X^{2}+153X+163$ at \code{n = 22};
\item the snippet ends with \code{set coefficient\_ring = "F\_ell"}, so a notebook session is left at the default.
\end{lookfor}

%======================================================================
\section{The ordinary rank in one call (\texttt{pow})}\label{sec:pg-pow}

For the ordinary rank of \S\ref{sec:pg-hida} one needs the rank of $U_\ell^{d}$, $d=\dim H^{1}$: the power $U^{d}$ annihilates the nilpotent
(non-ordinary) part and is invertible on the ordinary part, so $\mathrm{rank}\,U^{d}$ is the ordinary rank, and it equals $\mathrm{rank}\,U^{d+1}$.
The verb \code{pow(A, e)} (square-and-multiply, $A^{0}=I$) replaces the multiplication loop of \file{hida\_control\_13.oc}. At $\ell=13$ and
$n\equiv10\pmod{12}$ the dimension of $H^{1}$ grows by $2$ with each step of $12$ in $n$, while the ordinary rank stays $3$: $\Delta^{\pm}$ and the Eisenstein class
(Hida's control theorem, [U,~sec:u-lambda]; \CMPD).

\snippet{pow_ordinary_rank_13}

\begin{lookfor}
\item \code{dim H\^{}1} runs $3,5,7,9,11$ while \code{ordinary rank} is $3$ throughout; \code{rank U} is already $3$, i.e.\ $U_{13}$ itself
  vanishes on the non-ordinary part in these cells (a longer Jordan block there would make \code{rank U} exceed \code{rank(U\^{}d)}).
\end{lookfor}

%======================================================================
\section{The Hida tower of a weight class (\texttt{hida\_tower})}\label{sec:pg-hidatower}

The weight tower of {\ttfamily[C,~def:hidatower]} at the class $n\equiv10\pmod{12}$, $\ell=13$ is read in one call: the ordinary part of
$H^{1}(\SL_2(\Z);\Sym^{n}\,\cdot)$ over $\Z/13^{e}$ for $e=1,2$ on the reduction of the module over $\Z/13^{K}$ and on its free block
({\ttfamily[C,~cor:ordprecision]}), their agreement, and the characteristic polynomial of $T_2$ on the ordinary part. This is the computation of
\file{hida\_control\_13.oc} (\S\ref{sec:pg-hida}) as a verb (\CMPD).

\snippet{hida_tower_13_10}

\begin{lookfor}
\item \code{ord\_rank 3} at both levels, \code{readings agree: true}, \code{image\_free: true};
\item the level-$2$ polynomial $x^{3}+27x^{2}+75x+72$ reduces to the level-$1$ polynomial modulo $13$.
\end{lookfor}

%======================================================================
\section{Hida's control on forms (\texttt{hida\_lambda})}\label{sec:pg-hidalambda}

\code{hida\_lambda(ell, k, W)} computes the ordinary part $eS_k$ over $\Z/\ell^{W}$ and the characteristic polynomials of $T_p$ and $U_\ell$ on it
([U,~sec:u-lambda]). For $\ell=11$, $W=2$ the weights $12$ and $122=12+110$ are congruent modulo $(\ell-1)\ell^{W-1}=110$: although
$\dim S_{122}=9$ against $\dim S_{12}=1$, both ordinary parts have rank $1$ and the same polynomials for $T_2$ and $U_{11}$ over $\Z/11^{2}$, as Hida's control theorem
predicts (\CMPD).

\snippet{hida_lambda_11_12}

\begin{lookfor}
\item \code{dim S} $1$ and $9$, \code{rank e S} $1$ and $1$;
\item identical \code{charpoly T\_2} and \code{charpoly U} for $k=12$ and $k=122$.
\end{lookfor}

%----------------------------------------------------------------------
\section{The first ordinary level-one block of Loewy length 3: \texorpdfstring{$(484,547)$, $k=486$}{(484,547), k=486}}\label{sec:pg-ord547}

\paragraph{Ordinary rows have period $\ell-1$.} By Hida's control theorem {\ttfamily[C,~thm:hidacontrol]}, {\ttfamily[C,~thm:hidaweight]},
ordinary data depend on the weight only modulo $\ell-1$. The census shows this at the level of the full module. Ordinarity ($U_\ell$ a unit on
the block) is not a census column; it is read from the cells' records. In the signed census~4.0, $57\,728$ rows are ordinary ($27\,539$ main,
$30\,189$ b1); the $18\,124$ rows of the $580$ main cells without $U_\ell$ are left out. Pairing each cell $n\ge\ell$ with $n+\ell-1$ and comparing
the multiset of (residue degree, Loewy length) over the ordinary rows of $H^{1}$, the census is exactly periodic at every $\ell=17,\dots,199$:
$302$--$484$ pairs per prime, $0$ differing (\CMPD).
\begin{itemize}[leftmargin=1.6em,itemsep=1pt]
\item \emph{Loewy length.} Exactly one ordinary row has Loewy length $\ge3$: $(484,547)$, ideal~1, the exception below. (The investigation of
  2026-10-10 found none among the $16\,865$ ordinary rows of 3.4-pre.1, which did not yet contain $\ell=547$.)
\item \emph{The $240$ ordinary rows of Loewy length $2$} occur only at the primes $37$, $59$, $67$, $71$, $73$, $89$, $101$, $103$, $547$,
  $691$ (main) and $113$, $131$, $139$, $149$, $157$, $163$, $173$, $179$, $181$ (b1). $106$ have a reducible residual representation and $134$ an irreducible one. On 3.4-pre.1
  (then $96$ rows) the investigation identified the two mechanisms behind this split:
  \begin{itemize}[itemsep=0pt]
  \item an Eisenstein--cusp congruence at an irregular pair (the reducible rows; $56$ of the $96$);
  \item a cusp block over a ramified quadratic extension of $\Q_\ell$ (the irreducible rows; $40$ of the $96$, for instance $(66,89)$, $k=68$).
  \end{itemize}
  The new rows of~4.0 have not been examined one by one; that they follow the same two mechanisms is \EMP.
\end{itemize}

\paragraph{The exception $\ell=547$.} $547$ is irregular at $k\in\{270,486\}$, with $v_{547}(B_k)=1$ at both. At $k=486$ ($n=484$),
\emph{two} cusp eigenforms are congruent to $E_{486}$. They are conjugate over a ramified quadratic extension of $\Q_{547}$. The census cell
$(484,547)$ reads (\CMPD):
\begin{itemize}[leftmargin=1.6em,itemsep=1pt]
\item ideal~1 has $D=5=2g+1$;
\item $A=\F_{547}[x]/x^{3}$ (Hilbert function $(1,1,1)$, Gorenstein), and $H^{1}_\m\cong\F[x]/x^{3}\oplus\F[x]/x^{2}$;
\item the Eisenstein eigenvalues are $1+p^{485}\equiv147,539,3,377$ at $p=2,3,5,7$.
\end{itemize}

\paragraph{What the forms side and Hida theory add.} A forms-side computation with $T_{547}$ itself shows that the block is ordinary. With
Ohta's theorem and Hida control it also shows (\PROVED{} in the review cited below):
\begin{itemize}[leftmargin=1.6em,itemsep=1pt]
\item $\delta=0$;
\item the Loewy length is $3$ at every weight $k\equiv486\pmod{546}$.
\end{itemize}
The argument is modulo $\ell$ (the same ring $\T_{k,\m}/\ell$ and the same module at every such weight), so it does not depend on the
Eisenstein depth $v_{547}(B_k/k)$, which is $\ge2$ on one class modulo $546\cdot547$; it rests on Hida's control at the level of modules mod~$\ell$,
which the review cites from memory (not yet checked against the text).

\paragraph{How rare it is: a known exception.} That two cusp forms of weight $486$ are congruent to $E_{486}$ mod $547$ is
\textbf{known}: Bella\"iche and Pollack \cite[\S1.1, \S3.9]{ph_bp19} checked their hypothesis (Cuspidal rank one) for all irregular pairs
$(p,k)$ with $p<10^{5}$, and $(547,486)$ is the only exception, with cuspidal Eisenstein Hida--Hecke algebra of rank $2$. The project's own
screen of the $487$ irregular pairs with $\ell<8000$ (an Eisenstein generalised eigenspace of dimension $\ge3$ only at $\ell=547$, $k=486$) is
consistent with theirs and smaller. What this section adds, as far as the search saw, is the reading on cohomology: the ramified conjugacy, the
full ordinary block $\F[x]/x^{3}$, the module $\F[x]/x^{3}\oplus\F[x]/x^{2}$ with $\delta=0$ --- relevant to their Gorenstein hypotheses
\cite[\S1.2]{ph_bp19} --- and so the first ordinary level-one block of Loewy length $3$ ($g=2$ on the forms side: [EX,~ex:families]).

\noindent Sources: \pfile{docs/implementation/slot22/ordinary_loewy_investigation.md}, \pfile{ordinary_loewy_review.md},
\pfile{census34_lowl_ingest_report.md} \S5; the 4.0 counts by \file{docs/latex/tools/pg\_census\_check.py} (the ordinarity join of
\pfile{docs/implementation/slot22/ordloewy/ord_join.py}, on the census records of the project).

%----------------------------------------------------------------------
\section{Hida's control along \texorpdfstring{$n\equiv10\pmod{690}$}{n = 10 mod 690} at \texorpdfstring{$\ell=691$}{l=691}, to \texorpdfstring{$n=6910$}{n=6910}}\label{sec:pg-hi11}

\paragraph{The prediction.} The checkable {\ttfamily[C,~Hi11]} predicts the following along the Hida family of $\Delta$ glued to $E_{12}$:
\begin{itemize}[leftmargin=1.6em,itemsep=1pt]
\item the ordinary part of $H^{1}(\Ga;\Sym^{n}\F_{691}^{2})$ has rank $3$ at every $n\equiv10\pmod{690}$;
\item the characteristic polynomials there are $(X-(1+2^{11}))^{3}$ for $T_2$, $(X-(1+3^{11}))^{3}$ for $T_3$, and $(X-1)^{3}$ for
  $U_{691}$.
\end{itemize}
\S\ref{sec:pg-hida} is the same phenomenon at $\ell=13$.

\paragraph{The evidence.} Earlier runs held it at $n=700$, $1390$, $2080$, $2770$, $4150$ [EX,~nc:hi11]. The second-cluster job $3720162$
(2026-10-18) adds four cells, each with no $691$-torsion in $H^{1}$ ($e_{\max}=0$) and \code{verdict=HELD} (\CMPD):
\begin{center}\small
\begin{tabular}{@{}rrr@{}}
\toprule
$n$ & $\dim H^{1}$ & peak RSS (GiB) \\
\midrule
$4840$ & $807$ & $59.8$ \\
$5530$ & $923$ & $77.6$ \\
$6220$ & $1037$ & $97.9$ \\
$6910$ & $1153$ & $120.5$ \\
\bottomrule
\end{tabular}
\end{center}
So the ordinary rank stays $3$ while $\dim H^{1}$ grows past $1100$.

\paragraph{What is a theorem here.} The ordinary rank $3$ and the three characteristic polynomials along the family follow from Hida's
control theorem mod~$\ell$ at level one ($\ell\ge5$, $k\ge3$; \S\ref{sec:pg-hida}) together with the base weight $n=10$, where $\Delta$ is
$691$-ordinary and Ramanujan's congruence holds. [C] tags (Hi11) as a conjecture because it has not imported that theorem; the runs are a
check of it and of the probe. Not part of the theorem is ``no $691$-torsion'', and it fails once: $n=3460=5\cdot692$ is a Dickson degree at
$691$ (the only one in the family up to $n=6910$), so $H^{1}$ has $691$-torsion there; the probe refused that cell by name (job $3464192$,
$e_{\max}=1$) [EX,~nc:hi11], and the table above has the other members only.

\noindent Source: the ingest report of the second-cluster job $3720162$ (run data, out of tree).

%======================================================================
\part{Integral cohomology and torsion}\label{part:pg-integral}

Over $\Z$ the cohomology $H^{1}(\Ga;\Sym^{n}\Z^{2})$ is a finitely generated abelian group: a free part, whose rank is the count of
\S\ref{sec:pg-dims}, and torsion. The torsion lives at $2$ and $3$ (the orders of the finite stabilisers) and at the primes $\ell\ge5$ for which
$n$ is a Dickson degree, and it is exactly what made the mod-$\ell$ dimensions of \S\ref{sec:pg-dims} jump. The project computes these groups,
the integral Hecke operators on them and the boundary pair by the \emph{integral descent}: a fracture square that glues the cohomology over
$\Z[1/6]$ with the $2$- and $3$-adic pieces {\ttfamily[C,~sec:integral]}, {\ttfamily[C,~sec:descent]}, certified by
{\ttfamily[C,~thm:descentcert]} (the groups and the Hecke data) and {\ttfamily[C,~prop:descentpair]} (the boundary pair). This part shows the groups and their exponent law, the large-scale check of that law at $\ell=2,3$, the Hecke
operators at $\ell=2,3$, and how the Eisenstein congruence primes of Part~\ref{part:pg-congruences} reappear in the integral Hecke operators.

%======================================================================
\section{\texorpdfstring{$H^{1}$}{H1} over \texorpdfstring{$\Z$}{Z}: torsion primes and the exponent law}\label{sec:pg-descent-torsion}

The verb \code{descent(N, n)} (under \code{set coefficient\_ring = "Z"}; $N\le300$, $n\le60$ and $\psi(N)(n+1)\le300$, the Hecke primes
summing to at most $100$; larger cells by a campaign, [CLI,~sec:campaign-z]) returns the certified descent record of
$H^{q}(\Ga_0(N);\Sym^{n}\Z^{2})$, $q\le3$ (higher degrees are periodic) {\ttfamily[C,~thm:descentcert]}. Two statements make its torsion
predictable from $n$ alone.
\begin{itemize}[leftmargin=1.6em]
\item \emph{Which primes} (\PROVED, {\ttfamily[C,~prop:badprimes]}, with {\ttfamily[C,~cl:bockstein]} for the isomorphism; at level $N$,
  {\ttfamily[C,~prop:badprimesN]}). At level one and for $\ell\ge5$, $H^{1}(\Ga;\Sym^{n}\Z^{2})[\ell]\cong H^{0}(\Ga;\Sym^{n}\Fl^{2})$, so the
  $\ell$-primary torsion is non-zero exactly when $n$ is a Dickson degree at $\ell$, $n\in(\ell+1)\Z_{\ge0}+\ell(\ell-1)\Z_{\ge0}$, and it has as many cyclic
  summands as there are Dickson monomials of degree $n$; it lies in $H^{1}$ only. Every $\ell>n+1$ carries no torsion. The record lists these
  primes as \code{bad\_primes}, with the count.
\item \emph{Which exponent} (\LIT, Callegaro--Cohen--Salvetti \cite[Thm.~4.5]{ph_ccs13}). The largest $\ell$-power order in the torsion is
  $\beta(n,\ell)=\max\{j+1:n=\ell^{j}m,\ m\text{ a Dickson degree}\}$; the lower bound is {\ttfamily[C,~rmk:bocksteinexp]}(2). At $\ell=2$ every
  $m\ge2$ is a Dickson degree, at $\ell=3$ the Dickson degrees are $4\Z_{\ge0}+6\Z_{\ge0}$; so $\beta(n,2)=v_2(n)+1$ except at powers of $2$, and
  $\beta(n,3)=v_3(n)+1$ except at $n=2\cdot3^{j}$, where it is one less.
\end{itemize}
The snippet prints $H^{1}$ for seven values of $n$ and, next to it, both predictions computed in the script from $n$. The run results that
check the exponent law on $\approx3000$ cells up to $n=10^{4}$ are \S\ref{sec:pg-torsexp}.

\snippet{descent_torsion_beta}

\begin{lookfor}
\item \code{bad primes} equals the list of predicted primes on every line: $5,11$ at $n=12$ ($12=2\cdot(5+1)=11+1$), $7$ at $n=16$, $5,17,53$ at $n=54$;
\item the largest $2$-power: $\Z/2^{4}$ at $n=16=2^{4}$ ($v_2+1=5$ would overshoot), $\Z/2^{3}$ at $n=12$ and $60$;
\item the largest $3$-power: $\Z/3^{2}$ at $n=18=2\cdot3^{2}$ and $\Z/3^{3}$ at $n=54=2\cdot3^{3}$, one less than $v_3(n)+1$ --- the exception of $\beta$;
\item $n=60$: two $5$-primary summands, $\Z/5\oplus\Z/5^{2}$, because degree $60$ has two Dickson monomials at $5$, $P^{10}$ and $Q^{3}$ ($\deg P=6$, $\deg Q=20$), with exponents
  $v_5(\gcd(10,0))+1=2$ and $v_5(\gcd(0,3))+1=1$ \cite[Thm.~4.5]{ph_ccs13}; compare \code{dim H\^{}0} at $\ell=5$ in \S\ref{sec:pg-dims}.
\end{lookfor}

%----------------------------------------------------------------------
\section{The torsion exponent is the Frobenius bound (Callegaro--Cohen--Salvetti), checked at \texorpdfstring{$\ell=2$}{l=2} and \texorpdfstring{$\ell=3$}{l=3} (runs)}\label{sec:pg-torsexp}

Write the $\ell$-primary torsion of $H^{1}(\Ga;V_n)$ as $\bigoplus_t\Z/\ell^{a_t}$ and put $e_1=\max_t a_t$.

\paragraph{The bound $\beta$.} Call $m>0$ a \emph{Dickson degree} at $\ell$ if $H^{0}(\Ga;M_m)\ne0$. The ring of
$\SL_2(\Fl)$-invariants is generated in degrees $\ell+1$ and $\ell(\ell-1)$, and $\Ga\to\SL_2(\Fl)$ is onto, so the Dickson degrees are the positive
elements of $3\Z_{\ge0}+2\Z_{\ge0}$ at $\ell=2$ (every $m\ge2$) and of $4\Z_{\ge0}+6\Z_{\ge0}$ at $\ell=3$. Put
\[
  \beta(n,\ell):=\max\{\,j+1:\ n=\ell^{j}m,\ m\text{ a Dickson degree}\,\}.
\]
Since $1$ is not a Dickson degree at $2$ and $2$ is not one at $3$,
\[
  \beta(n,2)=v_2(n)+1-[n=2^{j}],\qquad \beta(n,3)=v_3(n)+1-[n=2\cdot3^{j}] .
\]
In words: $v_2(n)+1$ alone overshoots by one at the powers of $2$, and $v_3(n)+1$ alone at $n=2\cdot3^{j}$, where the best factorisation is
$n=3^{j-1}\cdot6$.

\paragraph{The theorem} (\LIT). Callegaro, Cohen and Salvetti \cite[Thm.~4.5]{ph_ccs13} (equivalently their Thm.~3.7(2)) determine the whole
$\ell$-primary torsion of $H^{1}(\SL_2(\Z);\Sym^{n}\Z^{2})$, for every prime $\ell$ and every $n>0$: each Dickson monomial $P^{k}Q^{h}$ of degree
$n$ (with $P$, $Q$ the invariants of degrees $\ell+1$ and $\ell(\ell-1)$) gives one summand $\Z/\ell^{v_\ell(\gcd(k,h))+1}$. Maximising over
$(k,h)$ gives exactly
\[
  e_1=\beta(n,\ell).
\]

\paragraph{What [C] proves.} The lower bound, at every prime: {\ttfamily[C,~rmk:bocksteinexp]}(2) shows that if $m$ is a Dickson degree, then
$H^{1}(\Ga;V_{\ell^{j}m})$ has an element of order $\ell^{j+1}$; so $e_1\ge\beta(n,\ell)$. The equality is used from the literature, not proved
in [C].

\paragraph{The data agree with it} (\CMPD, $0$ exceptions). $e_1=\beta(n,\ell)$ on every cell computed; the data are the L23 section of
census~4.0 (a pre-release addendum, \S\ref{sec:pg-runs}; table \pfile{deep1_l23_torsion40.tsv}):
\begin{itemize}[leftmargin=1.6em,itemsep=1pt]
\item the L23 records, even $n=4,\dots,700$: $349/349$ cells at $\ell=2$, $349/349$ at $\ell=3$;
\item the DEEP-1 run, even $n$: $2158/2158$ cells at $\ell=2$ ($n=600,\dots,10^{4}$; $2572$ computations, since $414$ cells were run in
  both modes \code{r3} and \code{full}), $128/128$ cells at $\ell=3$ ($n=600,\dots,3840$; $146$ computations).
\end{itemize}
Examples: $e_1=13$ at $(8192,2)$; $e_1=6$ at $(972,3)$; and $e_1=5$ at $(486,3)$, where $486=2\cdot3^{5}$. For $\ell\ge5$ the same equality
was seen on the $\Z/\ell^{8}$ probe of {\ttfamily[C,~rmk:bocksteinexp]}(3); the $\ell^{2}$-torsion of [EX,~ex:depth] is an instance of the same
theorem. The groups printed by the descent snippet of \S\ref{sec:pg-descent-torsion} agree with the theorem's multiset exactly, which supports
that the two $\Sym^{n}$ conventions match. So the runs are a large-scale check of \cite{ph_ccs13} (and of the code), not a new law.

\paragraph{The cells where $\beta$ and $v_3(n)+1$ differ.} The DEEP-1 grid at $\ell=3$ contains no $n=2\cdot3^{j}$, so on it $\beta(n,3)=v_3(n)+1$.
The table checks the exception itself, beyond that grid, with a control cell where the two agree:
\begin{center}\small
\begin{tabular}{@{}lll>{\raggedright\arraybackslash}p{0.2\linewidth}@{}}
\toprule
cell & $v_3(n)+1$ & $\beta(n,3)$ (\cite[Thm.~4.5]{ph_ccs13}) & observed $e_1$ \\
\midrule
$(1458=2\cdot3^{6},3)$ & $7$ & $6$ & $6$ \\
$(2916=4\cdot3^{6},3)$ (control) & $7$ & $7$ & $7$ \\
$(4374=2\cdot3^{7},3)$ & $8$ & $7$ & $7$ (sum rules resolved) \\
\bottomrule
\end{tabular}
\end{center}
At $(1458,3)$ the law $v_3(n)+1$ overshoots by one, exactly as $\beta$ says. Source: \pfile{docs/implementation/slot24/L23S14_report.md}.

\paragraph{Two companion observations at $\ell=2$} (\EMP). Let $c_p$ be the least elementary-divisor exponent of $T_p-(1+p)$ on
$\Lambda^{\mathrm{fr}}=H^{1}(\Ga;V_n\otimes\Z_2)/\mathrm{tors}$ (the free quotient, not localised: at $\ell=2$ there is one block).
\begin{itemize}[leftmargin=1.6em,itemsep=1pt]
\item \emph{Constant floors.} $c_p=\min(v_2(p^{n}-1),\nu_p)$ is \PROVED{} {\ttfamily[C,~lem:cprule]}, where $\nu_p$ is the least
  elementary-divisor exponent of $T_p-(1+p^{n+1})$. The data show $\nu_3,\nu_5,\nu_7$ constant, equal to $7$, $9$, $8$, so
  \[
    c_3=\min(v_2(n)+2,7),\qquad c_5=\min(v_2(n)+2,9),\qquad c_7=\min(v_2(n)+3,8)
  \]
  on all $2572$ DEEP-1 computations at even $n$ to $n=10^{4}$. The floor $\ge7$ at $p=3$ (and $\ge7$ at $p=5$) is Rustom's theorem \cite[Thm.~8.1, Prop.~8.4]{ph_rustom}, an
  eigenvalue congruence $a_p\equiv1+p^{k-1}\pmod{128}$ for every level-one eigenform; the sharpness and the values $9$, $8$ go beyond what
  our search found (2026-10-19).
\item \emph{Linear growth.} The largest finite elementary-divisor exponent of $T_\ell$ grows linearly, close to $(n+1)/(\ell+1)$. For
  example, it is $334$ at $(1000,2)$ and $249$ at $(1000,3)$.
\end{itemize}

\noindent Sources: the L23 section of census~4.0 (\pfile{deep1_l23_torsion40.tsv}, \pfile{deep1_l23_depths40.tsv},
\pfile{l23_torsion_law40.txt}; pre-release), recomputed by \file{pg\_census\_check.py}; \pfile{docs/implementation/slot24/L23S14_report.md}; \pfile{docs/implementation/slot22/census_l23_plan.md} \S4.2--4.3; \pfile{slot23/DEEP0_report.md} \S2;
\pfile{slot23/DEEP1_final_report.md}; \pfile{slot23/DEEP1_ell3_torsion_law.md}; the catalogue entry PH-01;
\pfile{docs/phenomenology/literature_search_B.md} \S1, \S3.

%----------------------------------------------------------------------
\section{\texorpdfstring{$\ell=2,3$}{l=2,3}: what the Hecke operators see (runs)}\label{sec:pg-l23}

At $\ell\in\{2,3\}$ the seven-vertex model computes \emph{Bredon} cohomology, not group cohomology, for $\Ga$ itself
({\ttfamily[C,~prop:modeltorsion]}: it computes group cohomology only for subgroups without elements of order divisible by $\ell$). The
statements below are about group cohomology: the L23 census computes it on the isotropy model $C_I$, cross-checked against the amalgam
resolution of $\Ga=\Z/4\ast_{\Z/2}\Z/6$ ({\ttfamily[C,~cl:amalgam]}), and the integral descent on its tree model.
\begin{enumerate}[leftmargin=1.8em,itemsep=2pt]
\item \PROVED{} \textbf{Prime-to-$\ell$ operators are scalars.} $T_p=(1+p)\cdot\mathrm{id}$ as a matrix on $H^{q}(\Ga;\Sym^{n}\Fl^{2})$ for
  every $p\ne\ell$ and every $q$ {\ttfamily[C,~prop:ltwothreescalar]}. So mod~$\ell$ the Hecke algebra is $\Fl[T_\ell]$, and the arithmetic is
  $\ell$-adic. The mechanism at level $\Ga(\ell)$ is known: Calegari--Emerton \cite[Lemma~3.1]{ph_ce_largeindex}, who credit Hatada
  (Math.\ Ann.\ 239 (1979), as quoted in \cite{ph_rustom}); the level-one form for every $q$ was not found. It is consistent with the large
  mod-$2$ Hecke algebra on forms (Nicolas--Serre) only through a large saturation defect: on the lattice the image of $\T/\ell$ is
  $\Fl[T_\ell]$, so $\dim(\T^{\mathrm{sat}}/\T)[\ell]=\dim M_k-2$. The L23 records show the hypothesis on the lattice ($T_p\equiv1+p$ and one unit
  divisor of $T_\ell$) on $349/349$ even cells at $\ell=2$ and at $\ell=3$ (the L23 section of~4.0, pre-release).
\item \EMP{} \textbf{$T_\ell$ is a rank-one idempotent.} On $H^{1}(\Ga;\Sym^{n}\Fl^{2})$, $n$ even, the characteristic polynomial of
  $T_\ell$ is $x^{h^{1}-1}(x-1)$. This holds on all DEEP-1 cells that compute it: $9303$ at $\ell=2$ ($n\le10^{4}$, odd $n$ included) and $196$ at $\ell=3$ ($n\le7500$). For
  example, at $(10\,000,2)$, $h^{1}=5001$, $\mathrm{rank}\,T_2=1$ and $\mathrm{rank}(T_2-1)=5000$. Only $\mathrm{rank}\,T_\ell\le2$ is proved
  {\ttfamily[C,~lem:tellrank]}. The eigenvalue shadow is Hatada's ($a_2\equiv0\bmod8$, $a_3\equiv0\bmod3$, via \cite{ph_rustom});
  for the matrix statement no prior was found in our search (2026-10-19).
\item \textsf{\small[LEAN]} \textbf{Closed dimensions.}
  \begin{itemize}[itemsep=0pt]
  \item $\dim_{\F_2}H^{q}(\Ga;\Sym^{n}\F_2^{2})=\lfloor n/2\rfloor+1$ for every $q\ge1$;
  \item $\dim_{\F_3}H^{q+2}(\Ga;\Sym^{n}\F_3^{2})=[n\equiv0,4\bmod 6]$ for even $n$ and $q\ge0$;
  \item at $\ell=3$, odd $n$, all groups vanish.
  \end{itemize}
  The rank formulas behind these values ($h^{q}=2d-a-s$ at $\ell=2$, and their analogues at $\ell=3$) are Lean theorems for every $n$; the
  values themselves are kernel-checked cell by cell for $n\le700$ (\textsf{\small[LEAN]} evidence for $n\le64$ and a few larger $n$, data beyond)
  {\ttfamily[C,~rmk:lvdimsltwothree]}. The closed forms above are read off those values, not formalised for every $n$. Odd $n$ at $\ell=2$ is
  {\ttfamily[C,~prop:oddnltwo]}.
\item \EMP{} \textbf{Cusp classes are divisible by $\ell$.} $H^{1}_{!}(\Ga;V_n)\subset\ell\cdot H^{1}(\Ga;V_n)$ on every computed cell,
  $n\le700$.
\item \EMP{} \textbf{A $\pm$-law at $\ell=3$.} For $n\equiv4\pmod6$, $T_3$ has the same elementary divisors on $\Lambda^{\mathrm{fr},+}_{!}$ and on
  $\Lambda^{\mathrm{fr},-}$, the $\iota$-eigenlattices of the free quotient over $\Z_3$ ($33/33$ cells, $n\le200$). There is no such law for $n\equiv0\pmod6$.
\item \CMPD{} \textbf{Integral examples.} The groups $H^{1}(\Ga;\Sym^{n}\Z^{2})$ themselves, with their $2$- and $3$-primary parts, are printed
  by the snippet of \S\ref{sec:pg-descent-torsion} (the integral descent, $n\le60$); for instance $H^{1}(\Ga;\Sym^{12}\Z^{2})\cong\Z\oplus\Z/2\oplus
  \Z/2^{2}\oplus\Z/2^{3}\oplus\Z/3\oplus\Z/3^{2}\oplus\Z/5\oplus\Z/11$ there. The number of $\ell$-primary cyclic summands is
  $\dim H^{0}(\Ga;M_n)$ {\ttfamily[C,~cl:bockstein]}; their exponents are the subject of \S\ref{sec:pg-torsexp}.
\end{enumerate}

\noindent Sources: the L23 section of census~4.0 (\pfile{docs/matrix-engine/results/congruence_census/edition3/4.0/README_l23.md}:
\pfile{deep1_l23_h01_40.tsv}, \pfile{deep1_l23_side40.tsv}, \pfile{l23_section40_v2.tsv}; pre-release);
\pfile{docs/implementation/slot22/census_l23_classification_review.md}, \pfile{L23_3_report.md}, \pfile{BIN3_report.md},
\pfile{slot23/DEEP1_final_report.md}.

%======================================================================
\section{Eisenstein congruence primes over \texorpdfstring{$\Z$}{Z}, and the boundary pair}\label{sec:pg-descent-congruence}

Ramanujan's congruence (\S\ref{sec:pg-ramanujan}) is invisible in the abelian group $H^{1}(\Ga;\Sym^{10}\Z^{2})=\Z^{3}\oplus\Z/2\oplus\Z/2^{2}\oplus\Z/3$: the
prime $691$ is not torsion. It is in the Hecke operators. On the free quotient, $T_2$ has the characteristic polynomial
$\chi_{\Z}(T_2)=(x-(1+2^{11}))(x-\tau(2))^{2}$ {\ttfamily[C,~prop:crtcharpolyN]}: one Eisenstein root, and the root of $\Delta$ twice (one class
in each $\iota$-eigenspace, \S\ref{sec:pg-sign}). Two roots of an integral characteristic polynomial are congruent modulo exactly the primes
dividing their difference, and $1+2^{11}-\tau(2)=2049+24=2073=3\cdot691$. So the Eisenstein congruence prime appears as a resultant of the
integral Hecke data, and its reduction is the non-semisimple block of \S\ref{sec:pg-ramanujan}. (The factor $3$ is the universal congruence at $\ell=3$, where every $T_p$ acts as $1+p$
(\S\ref{sec:pg-l23}).)

The same at level $11$ and weight $2$ ($n=0$), Mazur's original setting: $H^{1}(\Ga_0(11);\Z)=\Z^{3}$, the Eisenstein class with
$T_2=1+2=3$ and the elliptic curve $11a$ with $a_2=-2$; the difference is $5$, the numerator of $(N-1)/12=10/12$, Mazur's Eisenstein prime
(\LIT). Modulo $5$ the three roots collapse: \code{level\_n(11, 0, 5)} has $\chi(T_2)\equiv(x-3)^{3}$. The integral Hecke data are read on the
descent's certified groups {\ttfamily[C,~sec:descenthecke]}.

Last, the boundary pair over $\Z$ {\ttfamily[C,~sec:descentpair]}: the cohomology of $\Ga_0(11)$, of the boundary of its Borel--Serre
compactification (two cusps, so $H^{0}_{\partial}=\Z^{2}$), and the compactly supported cohomology $H^{*}_{c}$ as the cone of the restriction
--- the integral version of the sequence of \S\ref{sec:pg-sign}. $H^{2}_{c}=\Z$ is the fundamental class of the compactified modular curve, and
the rank-$2$ image of $H^{1}_{c}\to H^{1}$ is the interior cohomology of the genus-one curve $X_0(11)$.

\snippet{descent_congruence_primes}

\begin{lookfor}
\item the computed coefficients of $\chi_{\Z}(T_2)$ equal the \code{expected} expansion, at level $1$ and at level $11$;
\item \code{a - t = 2073 = 3 * 691} and \code{a - t = 5}: the congruence primes;
\item the mod-$5$ characteristic polynomial \code{[3, 2, 1, 1]} $=(x-3)^{3}$ (coefficients low to high);
\item the pair: \code{boundary H\^{}0 = Z\^{}2}, \code{compact H\^{}0 = 0}, \code{compact H\^{}2 = Z}, and torsion only at $2$ (the torsion primes at $n=0$
  are $2,3$, and here only $2$ occurs).
\end{lookfor}

%======================================================================
\section{When the torsion does not split off: non-split torsion in \texorpdfstring{$H^{1}$}{H1} over \texorpdfstring{$\Z$}{Z}, and what Galois sees (runs)}\label{sec:pg-nonsplit}

The integral descent of \S\ref{sec:pg-descent-torsion} writes $H:=H^{1}(\Ga_0(N);V)$, $V=\Sym^{n}\Z^{2}$, $n>0$, as $\mathcal T\oplus L$ with $\mathcal T=\operatorname{tors}H$ and $L=H/\mathcal T$
free. As an abelian group the sum is trivial. The Hecke operators need not respect it: in the basis torsion first, then free,
$T_p=\smat{\tau_p}{M_p}{0}{F_p}$ {\ttfamily[C,~prop:torsionhecke]}, and the question is whether the sequence
\[
  0\to\mathcal T_\ell\to H_\ell\to L_\ell\to0\tag{$\varepsilon$}
\]
splits as a sequence of modules over the Hecke algebra. This section is about the first examples where it does not, and about what that means. Tags: \PROVED{} (in [C] or proved here), \PGCOND{} (proved here under a named standard hypothesis), \CMPD{} (exact computation, command given), \LIT{} (cited from the literature), \EMP{} (a regularity seen in the stated range, not proved), \PGSPEC{} (a consistent reading, not proved).

\subsection*{1.\quad The test}

Write $\mathbb T'$ for the algebra of the $T_p$ with $p\nmid\ell N$, $c_p:=\tau_p|_{\mathcal T_\ell}$, $B_p:=F_p|_{L_\ell}$. A complement of $\mathcal T_\ell$ is the graph of some $h\in\Hom(L,\mathcal T_\ell)$, and it is $T_p$-stable iff $hB_p-c_ph=M_p$ {\ttfamily[C,~prop:heckeZglue]}(2)(b). No solution for the computed $T_p$ certifies that $\varepsilon$ is non-split for the whole Hecke algebra (the record's word \code{non\_split}). A solution is weaker: it gives a complement stable under $\mathbb T'_{\mathcal P}:=\Z[T_p: p\in\mathcal P]$ only, for the finite set $\mathcal P$ of primes computed (the field \code{P=[\dots]} of the \file{.out} files). The record's word for it is \code{split\_on\_P}, and {\ttfamily[C,~sec:recordschema]} says explicitly that this is not \code{split} for $\mathbb T$: that would need $\mathcal P$ to generate the image of $\mathbb T$ on $H^{1}_\ell/\ell^{K}$, torsion included ({\ttfamily[C,~rmk:heckeZcbthree]}). \emph{Convention for this section.} A computed verdict ``split'', written split$_{\mathcal P}$, always means \code{split\_on\_P}, and so does the bare word \code{split} printed by the tool; every ``exactly'' built on such a verdict is a statement on the computed primes. The word ``non-split'' is a certificate and needs no qualifier. A solution on a saturated $\mathbb T'$-stable sublattice restricts to any smaller such sublattice, so split$_{\mathcal P}$ passes to sublattices (\PROVED). The non-split test, again, does \emph{not} give an element of the glue $J$: split implies $J=0$, and the converse fails (2)(c). Nor does a non-zero $M_p$ matter by itself: $M_p$ changes with the basis, by $hB_p-c_ph$ ({\ttfamily[C,~sec:recordschema]}, item~8).

\subsection*{2.\quad What the torsion is, and how $T_p$ acts on it}

\begin{proposition*}[\PROVED{} here]\label{prop:pg-tors} For $n>0$ and $e\ge1$ put $M_e:=(V/\ell^{e})^{\Ga_0(N)}$. Then $\mathcal T_\ell=\varinjlim M_e$, and for $p\nmid N\ell$, $v\in M_e$,
\[ T_pv=(1+p)\,\sigma_p v,\qquad \sigma_a:=\rho\bigl(\mathrm{diag}(1,a)\bigr)\big|_{M_e},\quad \sigma_a\sigma_b=\sigma_{ab},\quad \sigma_a^{2}=a^{n}\qquad(a\in(\Z/\ell^{e})^\times). \]
\end{proposition*}
\begin{proof}
The first claim is the coefficient sequence for $0\to V\to V_\Q\to V_{\Q/\Z}\to0$ and $V_\Q^{\Ga}=0$, equivariantly by {\ttfamily[C,~lem:bocksteinallell]}'s argument. Let $g$ have integer entries, $\det g=p$, $c(g)\equiv0\bmod N$. In $\GL_2(\Z/\ell^{e})$ the element $\bar g\,\mathrm{diag}(1,p)^{-1}$ has determinant $1$ and $c\equiv0$ modulo the $\ell$-part of $N$; $\SL_2(\Z)\to\SL_2(\Z/N\ell^{e})$ is onto, so it is the image of some $\gamma\in\Ga$. Hence $\rho(g)v=\rho(\gamma)\rho(\mathrm{diag}(1,p))v=\sigma_pv$, $\sigma_pv$ being $\Ga$-invariant because $\mathrm{diag}(1,p)$ normalises the image of $\Ga$. Summing over the $1+p$ cosets of $\Ga\,\mathrm{diag}(1,p)\Ga$ gives the formula. Multiplicativity is clear. Since $\mathrm{diag}(a,1)=\mathrm{diag}(a,a^{-1})\mathrm{diag}(1,a)$ and $\mathrm{diag}(a,a^{-1})$ is in the image of $\Ga$, $\rho(\mathrm{diag}(a,1))$ acts as $\sigma_a$ on $M_e$; applying this to $\sigma_av$ gives $\sigma_a^{2}=\rho(aI)=a^{n}$.
\end{proof}
At level~$1$, mod~$\ell$, this is {\ttfamily[C,~prop:heckeH0]}(1). \emph{Example}: the vector $x^{n}$ lies in $M_e$ iff $a^{n}\equiv1\pmod{\ell^{e}}$ for all units $a$ and $\binom njc^{j}\equiv0\pmod{\ell^{e}}$ for $c\equiv0\bmod N$, $1\le j\le n$ (with [C]'s convention $\gamma\cdot x=ax+cy$); then $\sigma_ax^{n}=x^{n}$ and $T_p$ is exactly $1+p$ on the cyclic subgroup $\langle x^{n}\rangle\subseteq\mathcal T_\ell$ of order $\ell^{e}$. At $\ell\mid N$, $\ell$ odd, the largest such $e$ is $v_\ell(n)+1$ (for $\ell-1\mid n$) and at $\ell=2$ it is $v_2(n)+2$ ($n$ even) except that the term $j=1$ of the binomial condition intervenes, exactly when $\ell=2$ and $2\,\Vert\,N$, where it lowers the bound to $v_2(n)+1$: $e=4$ at $(8,4)$ and $(16,4)$, $3$ at $(8,10)$ and $(16,2)$, $1$ at $(27,2)$, $(27,4)$ (\PROVED; the table below shows it is the full exponent of $\mathcal T_\ell$ on the census cells, \CMPD).

\emph{The operators $\sigma_a$ need not be trivial, and $T_p$ does not notice.} At $(9,4)$, $\ell=3$, the space $M_1$ contains besides $x^{4}$ the vector $xy^{3}-x^{3}y$, and $\sigma_a(xy^{3}-x^{3}y)=\omega(a)(xy^{3}-x^{3}y)$ with $\omega$ the quadratic character modulo~$3$. On it $T_p=(1+p)\,\omega(p)$, which is $\equiv1+p\pmod3$ for every $p$, because $1+p\equiv0$ when $\omega(p)=-1$; so the flag \code{scal=True} of \file{nonsplit\_census.out} cannot detect it. Computing $M_e$ and $\sigma_a$ directly from $\Sym^{n}$, independently of the descent (mode \code{sigma}, the \code{SIGMA} rows of \file{nonsplit\_census.out}, whose orders agree with the torsion of the record on all $15$ lines computed), one finds $\sigma_a=1$ on all of $\mathcal T_\ell$ at the $\ell=2$ lines $(8,4)$, $(8,10)$, $(16,2)$, $(16,4)$, $(32,2)$, $(4,4)$, $(1,6)$ and at $(27,2)$, $(9,2)$, $(25,4)$; but $\mathcal T_3=\Z/3(1)\oplus\Z/3(\omega)$ at $(27,4)$, $(9,4)$, $(3,4)$, and $\mathcal T_3=(\Z/3\oplus\Z/9)(1)\oplus\Z/3(\omega)$ at $(9,6)$, $\mathcal T_3=(\Z/3)^{2}(1)\oplus(\Z/3)^{2}(\omega)$ at $(9,10)$ (\CMPD). Here $(\psi)$ means that $\sigma_a$ acts by $\psi(a)$ on the summand; \S5 returns to this.

\subsection*{3.\quad The census: where $\varepsilon$ is non-split}

HeckeZ edition~1 (\pfile{docs/integral_census/results/hecke_ed1/torsion_points.tsv}) has $12$ lines with \code{split=non\_split}, one for each of the anemic and the full algebra of the six cells below; at $\ell\mid N$ and none at level~$1$. The anemic algebra is reduced and its records are \emph{final}, except $(8,10)$, where two fields wait for an exact export word (\code{refused(export\_word)}, EXPORT-Z). The full algebra, which adds the $U_p$ at $p\mid N$, is non-reduced ($\code{redQ}=\code{non\_reduced\_certified}$) and its records were written \emph{pending} in edition~1, before package Z2 ($U_p$ over $\Z_{(\ell)}$); with the binary \code{cf93a542}, which has Z2, the six full records read \code{status=final}, generation \code{open(sandwich)} at each prime dividing $N$ with $\ell\mid N$ and C\textsubscript{T} \code{not\_run} (\code{C1=refused(non\_reduced)}), and the edition's rows are not re-issued. The verdict \code{non\_split} never waited: it uses only $T_p$ with $p\nmid N$, and it is the same on the anemic and on the full line of each cell. The columns $\mathrm{Eis}_1$, $\mathrm{Eis}_\chi$, $\mathrm{cusp}$ repeat the test on the saturated $\mathbb T'$-stable sublattices of $L$: the Eisenstein classes of trivial character, the twisted ones (eigenvalues $\chi(p)(1+p^{n+1})$, $\chi$ quadratic of conductor $4$ or $3$, $\mathrm{cond}^{2}\mid N$), and the cuspidal part; $S:=\mathcal T_\ell\cap\sum_p(T_p-1-p)H$ is explained in \S4. (\CMPD{} exact over $\Z$: \file{docs/latex/tools/pg\_nonsplit\_tors.py <binary> census} and \code{tbt}; output \file{phenomena\_gallery/nonsplit\_census.out}. The verdict on $L$ reproduces the record on every line. The columns other than $L$ are ours: the record states no verdict for them, and every entry of the table other than ``non-split'' is split$_{\mathcal P}$ in the sense of \S1. $\mathrm{Eis}_\chi$ is the Eisenstein part with $\chi$ of conductor~$4$ at $N=16$ and of conductor~$3$ at $N=9,27$.)
\begin{center}\small
\begin{tabular}{@{}llllllll@{}}
\toprule
$(N,n,\ell)$ & $\mathcal T_\ell$ & $\mathrm{rk}\,L$ & $L$ & $\mathrm{Eis}_1$ & $\mathrm{Eis}_\chi$ & cusp & $\lvert S\rvert$\\
\midrule
$(8,4,2)$ & $\Z/2,\Z/4,\Z/16$ & 10 & non-split & split$_{\mathcal P}$ & -- & non-split & 2 \\
$(8,10,2)$ & $(\Z/2)^3,\Z/4,(\Z/8)^2$ & 22 & non-split & split$_{\mathcal P}$ & -- & non-split & 2 \\
$(16,2,2)$ & $\Z/2,\Z/8$ & 12 & non-split & split$_{\mathcal P}$ & non-split & non-split & 2 \\
$(16,4,2)$ & $\Z/2,\Z/4,\Z/16$ & 20 & non-split & split$_{\mathcal P}$ & non-split & non-split & 4 \\
$(27,2,3)$ & $\Z/3$ & 18 & non-split & split$_{\mathcal P}$ & non-split & non-split & 3 \\
$(27,4,3)$ & $(\Z/3)^2$ & 30 & non-split & split$_{\mathcal P}$ & non-split & non-split & 3 \\
\bottomrule
\end{tabular}
\end{center}
Beyond the census. The family of mode \code{scan} consists of the prime-power levels and level $36$, $N\in\{2,3,4,5,7,8,9,16,25,27,32,36,49,64\}$ with $n\le12$ even, and the composite levels $N\in\{12,18,20,24,28,40,45,48,50,63\}$ with $n\le6$ even, in both cases subject to the console's bound $\psi(N)(n+1)\le300$: $78$ cells, by no means all cells under that bound (\file{phenomena\_gallery/nonsplit\_family.out}). It has $83$ lines with $\ell\mid N$. On them the verdict is non-split \emph{on the $35$ lines} with $\ell=2$, $8\mid N$, or $\ell$ odd, $\ell^{2}\mid N$, $(\ell-1)\mid n$, and on the other $48$ lines the test on $\mathcal P$ finds a complement, split$_{\mathcal P}$ (\CMPD; mode \code{law}, \code{LAW set=family}, $0$ violations). This is an \emph{empirical law} in this range, \EMP: $(N,\ell)=(4,2)$, $(12,2)$, $(20,2)$, $(28,2)$, $(3,3)$, $(5,5)$, $(7,7)$ and $(N,n,\ell)=(25,6,5)$ read split$_{\mathcal P}$. The clause $(\ell-1)\mid n$ rests on the single line $(25,6,5)$ (\code{scal=False}; the next test, $(49,6)$, is refused by the console's bound). Five further cells, not in the family and within the bound, confirm the law (\CMPD, \code{XCELL} and \code{XKZERO} rows of \file{nonsplit\_census.out}, \code{LAW set=extra}, $0$ violations): $(56,2)$ at $\ell=2$ (non-split, $K_0=2$ in the sense of \S7), $(8,14)$ (non-split, $K_0=2$), $(9,14)$ at $\ell=3$ (non-split, $K_0=1$), $(12,8)$ at $2$ and $3$ and $(5,16)$ at $5$ (split$_{\mathcal P}$). On every line of the family $\mathrm{Eis}_1$ reads split$_{\mathcal P}$; at $N=64$, and at $N=49$ on the line $\ell=3$ (where $p_0=2$), the computed lattice $\mathrm{Eis}_1$ contains twisted classes as well (a quadratic $\chi$ with $\chi(p_0)=1$ at the prime $p_0$ used to cut it); on the $\ell=2$ lines of $N=49$ the cut at $p_0=3$ is the trivial-character part, so the verdict there is the one for the larger lattice, and it passes to the trivial-character part by \S1 (\PROVED). The smallest non-split line is $(8,2,2)$, where $L$ is non-split while both $\mathrm{Eis}_1$ and the cuspidal lattice read split$_{\mathcal P}$: the non-splitting there lives in how $L_{\mathrm{Eis}}$ and $L_{\mathrm{cusp}}$ are glued inside $L$, not on either of them (\CMPD).

Note the pattern of the table: \emph{the trivial-character Eisenstein lattice reads split$_{\mathcal P}$ on every line}, and the twisted one is non-split on every line of the table where such a $\chi$ exists; $\mathrm{Eis}_\chi$ is absent at $N=8$ (no conductor~$4$ or $3$ character with square dividing~$8$) and those are the cells where the whole Eisenstein part is split$_{\mathcal P}$. (At $(8,2)$ this is not enough, as just said.)

\subsection*{4.\quad A canonical certificate: $S=\mathcal T_\ell\cap\mathbb I H$}

Assume that every $T_p$, $p\nmid\ell N$, acts on $\mathcal T_\ell$ by the scalar $1+p$; \code{scal=True} of \file{nonsplit\_census.out} records this, on every census line (it fails on some other cells, e.g.\ $(9,12)$, $\ell=3$, and at level one for $\ell\ge5$, and we then have no $S$ below). Let $\mathbb I$ be the ideal generated by the $T_p-1-p$. It kills $\mathcal T_\ell$, so $\mathbb IH\subseteq H$ maps onto $\mathbb IL$ with kernel $S:=\mathcal T_\ell\cap\mathbb IH$, a subgroup of $\mathcal T_\ell$, canonical once the scalar hypothesis holds. The tool computes it with the primes of $\mathcal P$; that is a subgroup $S_{\mathcal P}\subseteq S$, and so a non-zero value is still a certificate.

\begin{proposition*}[\PROVED{} here] If $\varepsilon$ splits for $\mathbb T'$ then $S=0$.\end{proposition*}
\begin{proof} If $H=\mathcal T_\ell\oplus L'$ with $L'$ stable, then $\mathbb IH=\mathbb IL'\subseteq L'$, so $\mathbb IH\cap\mathcal T_\ell=0$.\end{proof}
So a non-zero $S$ is a \emph{witness}: an explicit combination $\sum_p(T_p-1-p)x_p$ of classes that lands in the torsion and is not zero. On the census lines $S\ne0$ exactly where the verdict is non-split, on $\mathcal P$ (\CMPD; an exact form with $S$ and purity is \S7). On the six census lines $S$ is cyclic (the field \code{S\_inv}: $\Z/2$, $\Z/2$, $\Z/2$, $\Z/4$, $\Z/3$, $\Z/3$); in the family it is not always so, e.g.\ $(\Z/3)^{2}$ at $(9,10,3)$ and at the extra cell $(9,14,3)$. In the exported generators, one cyclic summand of $\mathcal T_\ell$ carries the whole class $\varepsilon$ (\CMPD). That this summand is the one generated by $x^{n}$ is what the exponents of \S2 suggest (\PGSPEC: the export's torsion generator is not compared with $x^{n}$).

\subsection*{5.\quad The Galois structure, and Eichler--Shimura}

\emph{Etale cohomology.} By Artin's comparison theorem (SGA~4, XI; cited from memory, not re-read for this section) applied to a finite-level cover with Hochschild--Serre,
$H_\ell\cong H^{1}_{\mathrm{et}}(Y_0(N)_{\bar\Q};\Sym^{n}R^{1}\pi_{*}\Z_\ell)$ as $\Z_\ell[T_p]$-modules, on the Deligne--Mumford stack $Y_0(N)$, so that $G_\Q$ acts on $H_\ell$, with $\mathcal T_\ell$ and $L_\ell$ stable and with $G_\Q$ commuting with every $T_p$, $p\nmid N\ell$. Everything below uses this comparison.

\begin{proposition*}[\PGCOND{} comparison] $G_\Q$ acts on $\mathcal T_\ell$ through the cyclotomic character: the geometric Frobenius at $p\nmid N\ell$ acts as $\sigma_p$.
\end{proposition*}
\begin{proof} On $M_e=H^{0}(\bar Y;\Sym^{n}R^{1}\pi_{*}\Z/\ell^{e})$, $G_\Q$ acts through $\pi_1(Y_\Q)$ modulo the geometric part. Its image in $\GL_2(\Z/\ell^{e})$ is the Borel-type group $\{c\equiv0\bmod\ell^{v_\ell(N)}\}$ with all determinants, and the geometric image is its determinant-one subgroup. Two elements of the same determinant differ by the geometric image, which fixes $M_e$; and $\sigma_p$ is the action of one with determinant $p$. \end{proof}

The torsion is therefore Galois-theoretically simple: it is unramified outside $\ell$, and $\sigma_p^{2}=p^{n}$. It need not be trivial: by \S2 the action on $\mathcal T_3$ at $(27,4)$ is $1\oplus\omega$. The essential input for the extension is the integral Eichler--Shimura relation.

\begin{center}\fbox{\parbox{0.92\linewidth}{\textbf{Hypothesis (ES$_{\Z}$).} On $H_\ell$, for $p\nmid N\ell$, with $F_p$ the geometric Frobenius: $T_p=F_p+p^{n+1}F_p^{-1}$. \\ \LIT{} Deligne proves this for $\Z_\ell$-coefficients on the interior cohomology of level-$n$ moduli schemes (\cite[Prop.~4.3, 4.4(iii), 4.8]{ph_deligne_bourbaki}: the correspondence $T_p$ on the normalisation is $F+I_p^{*}V$, with $F$ the geometric Frobenius and $FV=p^{n+1}$; the diamond operator $I_p^{*}$ is $1$ on the trivial-character part). The extension to the open curve and to the stack needs only the restriction $H^{1}(Y_{\bar\F_p},\mathcal F)\hookrightarrow H^{1}(U_{\bar\F_p},\mathcal F)$ to the ordinary locus (injective by Gysin, since a point has no $H^{1}$ of local cohomology) where the correspondence splits as a disjoint union, the argument being Deligne's on an open curve with $\ell\ne p$ (a tame sheaf on a smooth curve). We flag this extension as \emph{sketched, not a citation of a literal statement}; the integral form on open cohomology is hypothesis ES$_\Z$ here (\PGCOND{} below).}}\end{center}

Check on the torsion: there $F_p=\sigma_p$, and $\sigma_p+p^{n+1}\sigma_p^{-1}=\sigma_p+p\,\sigma_p=(1+p)\sigma_p=T_p$ by $\sigma_p^{2}=p^{n}$. (\PROVED{} this consistency; it does not test the sign: on a summand where $\sigma_p=\pm1$ one has $p^{n}\equiv1$ there and the identity holds for $\sigma_p^{-1}$ as well; $T_p=F_p+p^{n+1}F_p^{-1}$ is symmetric under $F_p\leftrightarrow p^{n+1}F_p^{-1}$, and the sign enters below only through the weight hypothesis of (1). The other common normalisation $T_p=\sum\rho(g^{*})$, $g^{*}=p\,g^{-1}$, gives the same formula for $T_p$, since $p^{n}\sigma_p^{-1}=\sigma_p$.)

\begin{theorem*}[\PGCOND{} ES$_{\Z}$] Assume \textup{ES$_{\Z}$}. If $\varepsilon$ is non-split for $\mathbb T'$, then $\varepsilon$ is non-split as a sequence of $\Z_\ell[G_\Q]$-modules. Its class is a non-zero element of $H^{1}(G_{\Q,N\ell\infty},\Hom(L_\ell,\mathcal T_\ell))$, with finite coefficients. Moreover, on a line $L'\subseteq L$ with $F_p=\lambda_p$ scalar, writing $F_p=\smat{\sigma_p}{b_p}{0}{\lambda_p}$ on $H'$, one has $T_p=\smat{(1+p)\sigma_p}{m_p}{0}{a_p}$ with
\[ m_p=\bigl(1-p^{n+1}\sigma_p^{-1}\lambda_p^{-1}\bigr)\circ b_p, \]
\end{theorem*}
\begin{proof} A $G_\Q$-stable complement is stable under $F_p$ and $F_p^{-1}$, hence by ES$_{\Z}$ under $T_p$: non-split for $\mathbb T'$ forces non-split for $G_\Q$. The formula is the $(1,2)$ entry of $T_p=F_p+p^{n+1}F_p^{-1}$, with $\sigma_p^{-1}$ acting on the values of~$b_p$.\end{proof}

So the Hecke module sees the Galois class only through the factor $\kappa_p:=p^{n+1}\sigma_p^{-1}\lambda_p^{-1}$, which depends only on the pair (summand of $\mathcal T_\ell$, line): $m_p=(1-\kappa_p)\circ b_p$. On a line with $\lambda_p=\chi(p)\,p^{n+1}$ and a summand on which $\sigma_a=\psi(a)$, $\psi$ of order at most~$2$ (which covers everything below), $\kappa_p=(\psi\chi)^{-1}(p)=\sigma_p\lambda_p^{-1}p^{n+1}$, and the Hecke algebra is \emph{blind} to $b$ on that pair exactly when $\psi\chi=1$. The cases are organised by this character.
\begin{enumerate}[leftmargin=1.8em,itemsep=2pt]
\item \PGCOND{} (ES$_{\Z}$ and the weight of the Eisenstein part: Frobenius on the Eisenstein part of $H^{1}$ is $p^{n+1}$ times a root of unity.) \emph{Trivial character on both sides, $\psi=1$, $\chi=1$}: $\kappa_p=1$, so $m_p=0$ for every $p$, whatever $b$ is: the pair is blind. This explains the column $\mathrm{Eis}_1=$ split$_{\mathcal P}$ \emph{on the summand of $\mathcal T_\ell$ where $\sigma=1$} (all of $\mathcal T_\ell$ on the lines of \S2 with $\sigma_a=1$; the $x^{n}$-line in general). It says nothing about the Galois class on these lines: there $b$ may be non-zero while $\mathbb T'$ splits.
\item \PGCOND{} \emph{Twisted line, $\sigma=1$ summand, $\psi=1$}: $\lambda_p=\chi(p)p^{n+1}$, where the Galois action on the (non-rational) cusps supplies the character $\chi$, so
\[ m_p=(1-\chi(p))\,b_p \qquad(\text{$\chi$ quadratic}), \]
 and $m_p=0$ whenever $\chi(p)=1$ (\CMPD, the \code{ELINE} rows of \file{nonsplit\_census.out}). (A cusp of denominator $d\mid N$ is defined over $\Q(\zeta_{\gcd(d,N/d)})$ \LIT; so non-rational cusps occur exactly when some $\gcd(d,N/d)>2$: $N=16$ and $N=9,27$ in the table, not $N=8$. Standard; not re-derived here. $N=25$ also has non-rational cusps and a twisted line that the tool does not isolate; no claim is made about it.)
\item \PGCOND{} \emph{The summand where $\sigma=\omega$, $\ell=3$} (it exists at $(27,4)$, $(9,4)$, $(9,6)$, $(3,4)$, $(9,10)$). On the trivial-character line $\psi\chi=\omega\neq1$: $m_p=(1-\omega(p))\,b_p$, \emph{not} blind, so the explanation in case~(1) does not apply here. There $\mathrm{Eis}_1=$ split$_{\mathcal P}$ is a computed fact: $M_p|_{\mathrm{Eis}_1}=0$ on the whole of $\mathcal T_3$ at all $17$ primes $p\le61$ prime to $3N$, at $(27,4)$, $(9,4)$, $(9,6)$ (\CMPD, the \code{EIS1M} rows). Under ES$_\Z$ this says that the $\omega$-component of $b$ on these lines, a homomorphism of $G_\Q$ into $\Z/3$ because $\omega\cdot\chi_{\mathrm{cyc}}^{-(n+1)}\bmod3$ is trivial, vanishes at the Frobenius of each such $p$ with $\omega(p)=-1$; if it did so at all $p$, it would vanish, by Chebotarev. On the \emph{twisted} line of those cells, with $\chi=\omega$ (quadratic of conductor~$3$), the pair (summand $\omega$, line) \emph{is} blind: $\psi\chi=1$. So $m_p$ on the twisted lines takes values in the $\sigma=1$ summand of $\mathcal T_3$, and this is the summand the verdict $\mathrm{Eis}_\chi$ non-split speaks about.
\item \CMPD{} \emph{The value of $b$ on the twisted lines, on the $\sigma=1$ summand.} On $(16,2)$, $(16,4)$, $(9,4)$, $(27,2)$, for the $17$ primes $p\le61$ prime to $N$ (line \code{ELINE\_CHECK}), $m_p=(1-\chi(p))\,t\,\log(p)$ on the one exported generator that carries the class, of order $\ell^{f}$, where $\log$ is the discrete logarithm ($\ell=3$: $\log p$ taken base $2$ in $(\Z/9)^{\times}/\{\pm1\}\cong\Z/3$; $\ell=2$: $\log p=k$ for $p\equiv\pm5^{k}$), and $t=2$ at $\ell=2$, $t$ a unit at $\ell=3$ (the scale of the line's generator). As a function of $p$ the right side depends only on $p$ modulo $2^{f}$ at $\ell=2$, modulo $3^{f+1}$ at $\ell=3$.\footnote{Computed only on the listed lines and primes $\le61$; at $\ell=2$ the identification of $b$ with a named class is \PGSPEC{} (see below).}
\end{enumerate}

\emph{What $b$ is on these lines.} $\Hom(L',\mathcal T_\ell)$ is not a trivial $G_\Q$-module: on the $\sigma=1$ summand $A$, of exponent $\ell^{f}$, the class lies in $H^{1}(G_{\Q,N\ell\infty},A(\psi_b))$ with $\psi_b(\mathrm{Frob}_p)=\sigma_p\lambda_p^{-1}=\chi(p)\,p^{-(n+1)}$, that is $\psi_b=\chi\,\chi_{\mathrm{cyc}}^{-(n+1)}$ (\PGCOND{}). At $\ell=3$, $f=1$, this is trivial ($\chi_{\mathrm{cyc}}\bmod3=\omega=\chi$ and $n$ is even), so $b\in\Hom(G_{\Q,\{3,\infty\}},\Z/3)$, a group spanned by the logarithm of the first layer of the cyclotomic $\Z_3$-extension (Kronecker--Weber, \LIT): at $(9,4)$, $(27,2)$ the non-zero $b$ is a unit multiple of it, which is the formula above. At $\ell=2$, $f\ge2$, the twist is not trivial: at $(16,2)$ it is $\psi_b(p)\equiv\chi_{-4}(p)\,p\pmod 8$, equal to $5$ at $p=3$. So $b$ is a cocycle, not a homomorphism of $\Z_2^{\times}$, and there is no statement to make inside $\Hom(\Z_2^{\times},\Z/2^{f})$. What the data say is weaker: the dependence on $p$ just stated, and at $(16,2)$, where $m_p\in\{0,4\}$: for the primes with $\chi(p)=-1$ the value $m_p/4\in\Z/2$ is the Kummer character of~$2$, that is $m_p=4$ exactly when $p\equiv3\pmod8$ (\CMPD, the \code{KUMMER2} row, on the $17$ primes).

\begin{proposition*}[\PGCOND{} ES$_{\Z}$, the weight hypothesis, and the computed verdict $\mathrm{Eis}_\chi=$ non-split] At $(16,2)$, $(9,4)$, $(27,2)$ the class $b$ of the twisted Eisenstein lattice is ramified at~$\ell$.\end{proposition*}
\begin{proof} Let $A$ be the cyclic summand carrying the class ($\Z/8$ at $(16,2)$; at $\ell=3$ the summand $\Z/3$ with $\sigma=1$, the other summand of $\mathcal T_3$ being blind by case~(3)), and $K$ the field cut out by $\psi_b$ on $A$: $K=\Q$ at $\ell=3$, $K=\Q(\sqrt2)$ at $(16,2)$ ($\psi_b$ is $1$ on $\pm1$ and $5$ on $\pm3$ modulo~$8$). Suppose the class of the $A$-component is unramified at~$\ell$. Since $N$ is a power of $\ell$, the restriction to $G_K$ is a homomorphism to the finite $\ell$-group $A$ unramified at every finite place of $K$, hence zero: $K$ has narrow class number~$1$ (\LIT; for $\Q(\sqrt2)$ the unit $1+\sqrt2$ has norm $-1$). So the class comes from $\Delta=\mathrm{Gal}(K/\Q)$. For $\ell=3$, $\Delta=1$ and it is zero; for $(16,2)$, $\Delta=\Z/2$ acts on $\Z/8$ by $5$, and $\widehat H^{1}(\Delta,\Z/8(5))=\ker(1+5)/(5-1)\Z/8=\{0,4\}/\{0,4\}=0$. Thus $b_p=(\sigma_p-\lambda_p)c$ for a fixed $c\in A$, and changing the complement by $c$ kills it, so $m_p=0$ on $A$. The rest of the torsion contributes nothing: at $(16,2)$ the summand $\Z/2$ has $m_p=(1-\chi(p))b_p$, even, so $0$; at $\ell=3$ the other summand is blind. So $\varepsilon$ would split on the lattice, against the certificate. \end{proof}
At $(16,4)$ the field $K$ is the real quartic subfield of $\Q(\zeta_{16})$ and the same argument would need Weber's parity theorem with signs; we have not written it out.

On the cuspidal lattice the theorem gives, \emph{on the $\sigma=1$ summand}, $M_p=b_p\circ(1+\rho_L(F_p)-T_p|_L)$ (use $p^{n+1}F^{-1}=T-F$ on $L$; on a summand with $\sigma=\psi$ one has $M_p=b_p-\psi(p)^{-1}b_p\circ(T_p-F_p)$), so again the Hecke class only sees $b$ through this factor. Beyond this we have no closed form for the class on $L_{\mathrm{cusp}}$; on the $\sigma=1$ summand the residual system is $T_p\equiv1+p$, i.e.\ $\bar\rho^{\mathrm{ss}}=1\oplus\omega$ (\PGCOND{} ES$_\Z$, with Brauer--Nesbitt on $L$), and the torsion carries only the direction ``$F_p=\sigma_p$'', not its companion~$p\sigma_p$. Beyond these two facts the cuspidal class is not identified.

\emph{What the class is not, and what is unknown.} \LIT{} The mod-$\ell$ phenomenon of Harder and Ribet (\S\ref{sec:pg-ramanujan}, \S\ref{sec:pg-kummer}) lives in the reduction of the free lattice and is tied to $L$-values: $691$ is not torsion at $(1,10)$, and these extension classes are not of that type. There is no theorem connecting $\varepsilon$ to a Selmer group $H^{1}_f$. A Ribet class is unramified at~$\ell$; the classes on the twisted Eisenstein lines above are \emph{ramified} at $\ell$ at $(16,2)$, $(9,4)$, $(27,2)$ (the Proposition above, \PGCOND{}, with its computed input), and they have finite coefficients. No motivic interpretation is known, and we do not imply one (\PGSPEC). At $\ell\mid N$ the curve has bad reduction at $\ell$, so no local condition at~$\ell$ is forced by what we compute. To be proved: (i) ES$_{\Z}$ in the form above on open integral cohomology; (ii) the value of $b$ on the twisted Eisenstein lines, conceptually; (iii) anything about the cuspidal class. \PGSPEC: the table suggests that the Eisenstein part of the class is controlled by the Galois action on the non-rational cusps; we have no proof.

\subsection*{6.\quad Torsion by torsion}

What should non-split \emph{torsion-by-torsion} mean? As in the torsion-by-free case, an extension that splits as abelian groups but not as Hecke modules; the natural reading is: $\mathcal T_\ell$ is a direct sum of $\mathbb T$-stable \emph{cyclic} groups, or it is not.

\begin{proposition*}[\PROVED{} here] Let $\ell$ be odd and $\mathbb T'$ the prime-to-$\ell N$ algebra. Then $\mathbb T'$ acts on $\mathcal T_\ell$ through scalars on the $\mu_{\ell-1}$-eigenspaces $M_{[\psi]}$: every abelian-group decomposition of $M_{[\psi]}$ is $\mathbb T'$-stable. Thus prime-to-$\ell N$ operators produce no torsion-by-torsion non-splitting at odd $\ell$.
\end{proposition*}
\begin{proof} $n$ is even (for odd $n$, $-I$ acts by $-1$, so $\mathcal T_\ell=0$ for odd $\ell$). For $a\in1+\ell\Z_\ell$ write $a=b^{2}$ with $b\in1+\ell\Z_\ell$; then $\sigma_a=\sigma_b^{2}=b^{n}=a^{n/2}$, a scalar. The prime-to-$\ell$ part of $\Z_\ell^\times$ acts semisimply, so $\mathcal T_\ell$ decomposes over its characters, and on each eigenspace $T_p=(1+p)\sigma_p$ is a scalar.\end{proof}
For $\ell=2$ the group $\Z_2^\times$ has a piece ($\pm1$, and $5$ not a square) the argument does not cover. Of the $136$ lines (cell, $\ell$) of the census cells, the extra cells $(4,4)$, $(9,4)$, $(9,6)$, $(3,4)$ and level $1$, $n\le60$, the $117$ with $\lvert\mathcal T_\ell\rvert\le1024$ were searched exhaustively and the other $19$ were not searched (\code{TBT\_SUMMARY}): at each of the $117$ the group $\mathcal T_\ell$ is a direct sum of $\mathbb T'$-stable cyclic groups (\CMPD). Also at $\ell=5$, $(1,60)$, where $\tau_p$ is not diagonal in the exported generators: it diagonalises by the change of basis $g_1\mapsto g_1+15g_2$.

The answer changes as soon as $T_\ell$ (or $U_\ell$) is in the algebra. \emph{There} torsion-by-torsion non-splitting exists, and it is {\ttfamily[C,~prop:heckeH0]}(3),(4) in the integral setting: \CMPD{} at level $1$, $\ell=2$, $n=6$ the torsion is $\Z/2\oplus\Z/4$ with $T_2=\smat{0}{0}{2}{2}$ (rows modulo the generators' orders) and every $T_p$, $p$ odd, the scalar $1+p$, and the module is not a direct sum of $T_2$-stable cyclic groups (exhaustive). Of the $117$ lines searched, $15$ are non-decomposable for the algebra with $T_\ell$ (\code{TBT} rows, \file{nonsplit\_census.out}): $\ell=2$ at $(1,6)$, $(1,10)$, $\dots$, $\ell=3$ at $(1,24)$, $(1,42)$, and $(9,6)$ at $\ell=2$; exhaustive search under the bound (\CMPD). The first $\ell\ge5$ instance is $(120,5)$, by {\ttfamily[C,~prop:heckeH0]}(4), outside the console's limits. $T_\ell$ is not in the image of $\Z_\ell[G_\Q]$, so Eichler--Shimura does not transfer this kind of non-splitting to Galois; and since $G_\Q$ acts on $\mathcal T_\ell$ through the operators $\sigma_a$ of \S2 ($a\in\Z_\ell^{\times}$), the proof of the Proposition shows that \emph{at odd~$\ell$ the Galois action on $\mathcal T_\ell$ is by scalars on each $\mu_{\ell-1}$-eigenspace} (for instance the characters $1$ and $\omega$ at $(27,4)$): no non-split torsion-by-torsion comes from Galois (\PGCOND{} comparison). At $\ell=2$ this is open.

Between degrees one can also ask about $0\to H^{1}(V)/\ell\to H^{1}(V/\ell)\to H^{2}(V)[\ell]\to0$ of Hecke modules (at $\ell\in\{2,3\}$ the two ends are non-zero, {\ttfamily[C,~sec:descent]}). The shipped binary does not expose $T_p$ on $H^{2}$ or on $H^{1}(V/\ell)$ with a basis compatible with this sequence; we did not test it.

\subsection*{7.\quad What the computations modulo $\ell$ see}

The extension is a statement about $\Z$-modules. The owner asked whether one can find it from the much cheaper $\F_\ell$-computations of the census. Put $H_K:=H\otimes\Z/\ell^{K}$, $A_K:=\mathcal T_\ell/\ell^{K}$ and let
$\varepsilon_K\colon0\to A_K\to H_K\to L/\ell^{K}\to0$, exact because $L$ is free.

\begin{proposition*}[\PROVED{} here; (2) uses {\ttfamily[C,~prop:heckeZglue]}(2)(a)]
\textup{(1)} If $\varepsilon_{K+1}$ splits for $\mathbb T'$ so does $\varepsilon_K$; so non-split modulo $\ell$ implies non-split over $\Z$. \textup{(2)} For $K\ge e_\ell$, $\varepsilon_K$ splits iff $\varepsilon$ does: if the exponent of $\mathcal T_\ell$ is $e_\ell=1$, the test modulo~$\ell$ is exact. \textup{(3) Exact test.} Suppose every $T_p$, $p\nmid\ell N$, acts on $A_K$ by a scalar $c_p$, and let $I_K:=\sum_p(T_p-c_p)H_K$. Then $\varepsilon_K$ splits iff $A_K\to H_K/I_K$ is injective with image a direct summand of the group $H_K/I_K$. For $K=1$: iff $A_1\cap I_1=0$.
\end{proposition*}
\begin{proof} (1) A complement reduces. (3) A complement $C$ contains $I_K$ (the $T_p$ act by $c_p$ on $H_K/C\cong A_K$), so $H_K/I_K=A_K\oplus C/I_K$. Conversely, the $T_p$ act on $H_K/I_K$ by the scalars $c_p$; so the preimage of an abelian-group complement of $A_K$ in $H_K/I_K$ is $\mathbb T'$-stable, and it is a complement of $A_K$ in $H_K$.\end{proof}

So $\F_\ell$ sees everything at $e_\ell=1$, and the counterexamples for $e_\ell>1$ exist. Let $K_0$ be the least $K$ with $\varepsilon_K$ non-split. On the $35$ non-split lines of the family of \S3 (\CMPD, the \code{KZERO} rows; all $35$ are non-split over $\Z$ by the exact test of \file{pg\_nonsplit\_tors.py}, and \code{hecke\_structure} reads \code{non\_split} in \code{G6} at the cells we also ran through it: $(8,2)$, $(9,2)$, $(9,4)$, $(25,4)$, $(25,8)$, $(32,2)$, $(36,2)$, $(64,2)$), $K_0=1$ except where the $2$-part of $N$ is exactly~$8$: $K_0=2$ at $(8,4),\dots,(8,12)$, $(24,2)$, $(24,4)$, $(40,2)$, and $K_0=3$ at $(8,2)$. \emph{At these levels the class dies modulo $2$}: $\varepsilon_1$ reads split$_{\mathcal P}$ while $\varepsilon$ is non-split, a counterexample to the converse of (1) for $\mathbb T'_{\mathcal P}$, which is what is computed. The same verdicts are read from the shipped mod-$\ell^{K}$ route \code{D.module(ell, K)}, independently of \code{module\_Z}, at $(8,4)$, $(8,2)$, $(16,2)$, $(27,2)$, $(25,4)$, and the control $(4,4)$ (\code{MODK} rows). On the lines where the scalar hypothesis holds, $\dim(A_1\cap I_1)>0$ exactly where $K_0=1$ (on $\mathcal P$).

\emph{Three warnings about what ``modulo $\ell$'' means.} (a) $H^{1}(\Ga;V/\ell)$ is larger than $H_1$: the universal-coefficient sequence $0\to H/\ell\to H^{1}(V/\ell)\to H^{2}(V)[\ell]\to0$ has a non-zero right term at $\ell=2$ (e.g.\ $H^{2}(\Ga_0(16);\Sym^{2}\Z^{2})=(\Z/2)^{2}$, printed by \code{descent(16, 2)}); it vanishes for $\ell\ge5$ and for $\ell=3$ when $\Ga_0(N)$ has no element of order $3$ (e.g.\ $9\mid N$), because then $H^{2}$ has no $\ell$-part ({\ttfamily[C,~prop:badprimesN]}; the examples $9$, $16$, $27$ happen to have no element of order $2$ in $\Ga_0(N)/\pm1$ either). (b) In these calculations $A_1$ is the image of the torsion, which a pure Hecke computation over $\F_\ell$ does not mark: without \code{D.module} one knows only the block. (c) The verb \code{level\_n(N, n, ell)} ($\ell\ge5$) prints characteristic polynomials: it cannot see any of this.

\emph{What a block alone can decide.} If $\mathfrak mM=0$ on an $\F_\ell$-block $M$ containing the torsion, then $M$ is a module over the field $\mathbb T/\mathfrak m$, so the image of $A_1$ in $M$ is a direct summand and the block splits, for $\mathbb T$ and a fortiori for $\mathbb T'$; no scalar hypothesis is used. This is the one place in the section where a split is certified, and it is not a \code{split\_on\_P}. Supports do the rest: {\ttfamily[C,~prop:heckeZsplit]}. (The census's $H^{0}$-block is $\mathcal T_\ell[\ell]$, while $A_1=\mathcal T_\ell/\ell$. They differ as $\mathbb T$-modules in general, but multiplication by $\ell$ on $\mathcal T_\ell$ shows they have the same composition factors, hence the same support, which is all the search below uses.) This gives a \emph{search} of the census, by \file{pg\_nonsplit\_tors.py modl}: every row of the signed census~4.0 (main and b1; level $1$, $\ell\ge5$; $448\,978$ rows, $n\le1000$, $46$ primes) whose $H^{1}$-ideal is matched by an $H^{0}$-ideal. There are $3236$ such collisions, and all of them have Loewy length~$1$ ($\mathfrak mM=0$; in $2718$ of them $L/\ell$ is non-zero on the block, the rows with a \code{sturm} object, field \code{free\_on\_block}). Of the $3236$, $3080$ are certified and $156$ are heuristic in the census's sense (\code{certified=heuristic}, \code{final=false}, at $\ell=17,19,23,29,31$ ($53$, $51$, $38$, $11$, $3$ rows), $n\ge552$; by {\ttfamily[E,~sec:census]}, column rule (d7), their module columns, the Loewy length included, are heuristic). So no level-one torsion point in the census range is non-split \emph{modulo~$\ell$}: $\varepsilon_1$ splits on the $3236$ rows, certified on $3080$ (\CMPD, from the census's data; the \code{MODL\_CERT} row). By (2), at the $2952$ rows with $\beta(n,\ell)=1$ (the exponent, \cite[Thm.~4.5]{ph_ccs13}) the extension $\varepsilon$ itself is split over $\Z$ ($2801$ certified, $151$ heuristic); at the remaining $284$ rows (152 cells, $\beta>1$) only $\varepsilon_1$ is split, and the census does not decide~$\varepsilon$. At $\ell\in\{2,3\}$, level $1$, the prime-to-$\ell$ operators are scalars on the whole of $H^{1}(\F_\ell)$ {\ttfamily[C,~prop:ltwothreescalar]}, so $\varepsilon_1$ splits for $\mathbb T'$ at every $n$ (\PROVED). HeckeZ edition~1, which tests over $\Z$ on $\mathbb T'$ to $n\le268$, reads \code{split\_on\_P} at every level-one point it tests, consistent with this.

For \emph{torsion by torsion} the same census row gives a criterion directly: if the block of $H^{0}(\Ga;V/\ell)\cong\mathcal T_\ell[\ell]$ is not semisimple for $\mathbb T$ (Loewy length $\ge2$), then $\mathcal T_\ell$ is not a direct sum of $\mathbb T$-stable cyclic groups (\PROVED: a $\mathbb T$-stable cyclic decomposition would restrict to one of $\mathcal T_\ell[\ell]$). Of the $3236$ $H^{0}$-rows at $n>0$ (the census has $5$ more $H^{0}$-rows, at $n=0$, where $H^{0}=\Z$ is not torsion), $30$ are non-semisimple, at level $1$, $\ell\ge5$: the cells $(n,\ell)=(120,5),(220,5),(240,5),(320,5),(336,7),\dots,(980,5)$ (\code{MODL\_H0} row; the first cells are the $\ell^{3}-\ell$ of {\ttfamily[C,~prop:heckeH0]}(4)). At $e_\ell=1$ this reading is exact, in general one-way. None of these $n$ is within reach of the descent over $\Z$ (the console's $n\le60$), and $T_\ell$ does not come from $G_\Q$. For the prime-to-$\ell N$ algebra (\S6) there is nothing to search: $H^{0}$ carries $(1+p)p^{a}$, semisimple.

In short: the prime-to-$\ell N$ torsion-by-free obstruction is \emph{exactly} visible modulo~$\ell$ when $e_\ell=1$; modulo $\ell$ it is sufficient but not necessary (level~$8$), and it needs the torsion to be marked: the census alone can only exclude it (Loewy length~$1$), not find it. At level~$1$ in the census range no torsion point is non-split modulo~$\ell$ (certified on $3080$ of the $3236$ rows); at $284$ rows with $\beta>1$ the extension over~$\Z$ stays undecided. At level $N$ one needs \code{D.module} at $K=e_\ell$ (the console's limits), the only route that marks the torsion. At level $N$ the law of \S3 (with $\ell=2$ the level must be divisible by $8$; with $\ell$ odd, by $\ell^{2}$, and $\ell-1\mid n$) is the empirical statement to test next, on $\mathcal P$; it is not proved, and any proof must explain why $\ell^{2}\mid N$ (resp.\ $8\mid N$) is what matters.

\emph{An open direction for the odd-$\ell$ half of the law} (\PGSPEC). For $\ell$ odd, $\ell^{2}\mid N$ is exactly the condition for a cusp with $\ell\mid\gcd(d,N/d)$, hence for the twisted Eisenstein line $E_{n+2}^{\chi_\ell,\chi_\ell}$ with $\chi_\ell$ quadratic of conductor~$\ell$ in $L$; and $(\ell-1)\mid n$ is exactly the condition $x^{n}\in M_1$, the $\sigma=1$ summand of \S2. Under ES$_\Z$ one has $m_p=(1-\chi_\ell(p))\,b_p$ on that pair, and $b$ restricted to $G_{\Q(\zeta_\ell)}$ is a homomorphism to $\Z/\ell$ unramified outside $\ell$, so it is the cyclotomic logarithm up to a unit if it is not zero. The odd-$\ell$ law would then reduce to the non-vanishing of one coefficient of the Eisenstein cocycle of $E^{\chi_\ell,\chi_\ell}$ on $x^{n}$, a Stevens-type computation possibly within reach for $\ell=3$. At $\ell=2$ the mechanism must differ ($N=8$ has only rational cusps and the non-splitting is cuspidal), so ``$8\mid N$'' needs another idea.

\subsection*{8.\quad Run it}

\snippet{nonsplit_torsion_16_2}

\begin{lookfor}
\item \code{status: final}, \code{gen=[2:open(sandwich)]}, \code{C\_T=not\_run} (this record is non-reduced, \code{C1=refused(non\_reduced)}; the transcript is that of the binary with sha256 prefix \code{cf93a542}, which has $U_2$, $U_3$ over $\Z_{(2)},\Z_{(3)}$), and \code{G6=[2.pt[0]=non\_split]}, the verdict, does not depend on that;
\item $T_3$, $T_5$, $T_7$ act on the $\Z/2$ and $\Z/8$ torsion generators as $0$, $4$; $0$, $6$; $0$, $0$: the scalars $1+p\bmod d_i$;
\item the row $M$ of the $\Z/8$ generator is not zero on the free basis (but $M$ alone proves nothing, \S1);
\item the verdict, the sublattice columns, $S$ and the torsion-by-torsion lines come from \file{docs/latex/tools/pg\_nonsplit\_tors.py}, not from this snippet.
\end{lookfor}

%======================================================================
\appendix
%======================================================================
\section{What the console cannot show yet, and reading \texorpdfstring{$\Sigma(691)$}{Sigma(691)}}\label{sec:pg-missing}

Every phenomenon of this gallery has a \code{read\_only} verb, shown once in \S\ref{sec:pg-pow}--\ref{sec:pg-hidalambda}; the one exception is
the $\Sigma(691)$ atlas store, which is read outside the console (next paragraph).
\begin{center}\small
\begin{tabular}{@{}>{\raggedright}p{0.42\linewidth}>{\raggedright\arraybackslash}p{0.52\linewidth}@{}}
\toprule
phenomenon & verb or reader \\
\midrule
level $N$, old forms & \code{level\_n(N, n, ell)}, \S\ref{sec:pg-leveln} \\
nebentypus characters & \code{nebentypus(N, m, n, ell)} ($m$ = Conrey label), \S\ref{sec:pg-nebentypus} \\
$p$-adic $L$-values, the measure congruence [E,~sec:lvalue] & \code{lfunction} \\
the $\Lambda$-adic Hecke algebra [U,~sec:u-lambda] & \code{hida\_lambda(ell, k, W)}, \S\ref{sec:pg-hidalambda} \\
Hida towers & \code{hida\_tower}, \S\ref{sec:pg-hidatower} \\
Kummer generators [E,~sec:lvnf] & \code{kummer\_unit}, \S\ref{sec:pg-kummerunit} \\
the forms side & \code{forms\_algebra(k, ell, primes, terms)}, \S\ref{sec:pg-forms} \\
a matrix power & \code{pow}, \S\ref{sec:pg-pow} \\
the $\Sigma(691)$ atlas store & not a verb: the notebook's atlas channel, or \code{hacs --atlas --query <store.json> '<request>'} (below) \\
\bottomrule
\end{tabular}
\end{center}
\paragraph{Reading $\Sigma(691)$.} A path from the user is file I/O, so no \code{read\_only} verb opens a store and \code{--serve} does not grow one. Two readers exist,
both on the same request handler (\code{info}, \code{cells}, \code{cell}, \code{system}, \code{systems}; the store is opened read-only, its FNV-1a-64 is in every reply, and the
path comes from the command line only): the notebook's atlas channel (\code{--serve-atlas}, protocol \code{oc-serve-atlas/1}, behind the gateway), and the one-shot
\begin{Verbatim}[fontsize=\small]
hacs --atlas --query <atlas dir>/ell691/hecke_atlas_ell691.json \
    '{"op":"cell","id":"c1","n":690}'
\end{Verbatim}
which prints one \code{atlas\_reply} frame on standard output (exit $0$ for \code{ok}, $3$ for an \code{error} or \code{rejected} reply, $1$ for a usage error or a refused store) and is
built only with \code{OC\_BUILD\_NOTEBOOK=ON}. On the real store of $\Sigma(691)$ (FNV \code{0x73c7210047a81b97}) a query takes about a second and $0.5$\,GB, dominated by the
open (the scan of the $452$\,MB file); several queries should share one \code{--serve-atlas} session. The reply of \code{cell} with \code{n = 690} names $13$ ideals and the dimension $115$ of $H^{1}$.

A per-block eigenvalue reader over $\Z/\ell^{K}$ (the ordinary characteristic polynomial of \S\ref{sec:pg-hida} in one call) is still a candidate; \code{hida\_tower} gives it
for level one.

%======================================================================
\section{Version history}\label{app:pg-history}

\begin{itemize}[leftmargin=1.6em,itemsep=2pt]
\item Version 1.5.4 (2026-11-09, release audit fixes): the program is called \file{hacs} in the snippets (\file{cohomology\_app} remains a link to it); the build instructions are cited as \file{BUILDING.md}; a source report whose path named a machine is described instead, and a cluster kit's path is dropped. No change of content.
\item Version 1.5.3 (slot~26, 2026-11-09): machine names and private paths removed for the public release (owner ruling 2026-11-09); no change of content.
\item Version 1.5.2 (slot~26, PG-NONSPLIT, 2026-11-09): referee (1.5.0) and confirm-read (1.5.1, \file{docs/implementation/slot26/PG\_nonsplit\_confirm.md}) passed, with two one-sentence repairs applied and nothing else changed: the sign check after hypothesis ES$_\Z$ no longer claims to test the sign ($F_p=\sigma_p^{-1}$ satisfies the identity too), and the twisted classes in the computed $\mathrm{Eis}_1$ occur at $N=64$ and at $N=49$, $\ell=3$ only. The status text about the non-reduced full records, written before package Z2 merged, now agrees with the binary \code{cf93a542} (status \code{final}, generation \code{open(sandwich)}, C\textsubscript{T} \code{not\_run}); snippet \file{nonsplit\_torsion\_16\_2} regenerated; no mathematics changed.\par\item Version 1.5.1 (slot~26, PG-NONSPLIT referee fixes, 2026-11-09): \S\ref{sec:pg-nonsplit} revised on the report of \file{docs/implementation/slot26/PG\_nonsplit\_referee.md}: computed splits are \code{split\_on\_P} on the printed $\mathcal P$; the operators $\sigma_a$ on $\mathcal T_\ell$ are computed and need not be trivial, the Eisenstein discussion is organised by the character $\psi\chi$; the coefficients of the class on the twisted lines and a proof of ramification under ES$_\Z$; corrected counts (census certification, $H^{0}$-rows at $n>0$, searched torsion-by-torsion lines, $S$ not always cyclic). The tool gains the modes \code{sigma}, \code{extra}, \code{law} and the field \code{S\_inv}.\par\item Version 1.5.0 (slot~26, PG-NONSPLIT, 2026-11-08; owner request; not yet refereed): new \S\ref{sec:pg-nonsplit} (torsion against the free part, and torsion by torsion, over $\Z$, with a Galois reading conditional on an integral Eichler--Shimura hypothesis stated there, and the question of what the computations modulo~$\ell$ see, with the exact test and the census search of its \S7), new snippet \file{nonsplit\_torsion\_16\_2} and tool \file{docs/latex/tools/pg\_nonsplit\_tors.py}. Regenerated outputs include the two snippet edits of da84747c (\file{descent\_congruence\_primes}, \file{hida\_control\_13}: comment lines only).
\item Version 1.4.3 (slot~25, 2026-11-06; LaTeX polish stage 1): uniform author block; the abstract no longer quotes figures that change with every edition. Removed
  verbatim from the abstract, item~III: ``and what census edition~4.0,
  $448\,978$ local Hecke algebras, found;'' and, from its last sentence, ``(2026-10-19)'' (``where a literature search (2026-10-19) found a source, it is
  cited''). Former author line: ``The orbifold-cohomology-of-$\GL_2$ project''. The figures remain in \S\ref{sec:pg-runs}. No number changed.
\item Version 1.4.2 (slot~25, 2026-11-06): two parenthetical notes, in \S\ref{sec:pg-census34} and \S\ref{sec:pg-dihedral3}: the $\ell=59$ candidate is
  verified (\pfile{docs/implementation/slot25/R59_verification.md}) and [E] 1.25.0-pre.1 restates (R) with ``of one weight'' for ``not dihedral'', by owner
  ruling of 2026-11-06. No number changed.
\item Version 1.4.1 (slot~25, PG141, 2026-11-06; referee \pfile{docs/implementation/slot25/PG140_referee.md}, R1--R6). R1: the rows at $\ell=59$ and $139$
  are no longer ``not checked'': $139$ has a congruence in weight $138$ ((R) silent), $59$ is congruence-free and not dihedral, a candidate counterexample
  to [E,~pr:latticesplit] (R) pending the owner (four-prime evidence). R2: the algebra $\F_q[x,y]/(x^{3},xy,y^{2})$ named. R3: the module route reads $A$, not $\mathsf A_1$, on
  a torsion block. R4: the hidden flag cited to [LG,~def:hiddencong], a notation concordance. R5: the $459$ \code{not\_stable} rows. R6: \file{pg\_census\_check.py}
  reads the printed literals from this file. No number changed.
\item Version 1.4.0 (slot~25, PG140, 2026-11-06; owner request of 2026-11-06; report \pfile{docs/implementation/slot25/PG140_report.md}).
  \emph{Notation:} the full localisation $\mathbb L_\m=H^{1}(\Ga;V_n\otimes\Zl)_\m$ (torsion included) and its free quotient
  $\Lambda^{\mathrm{fr}}_\m$ are named apart (\S\ref{sec:pg-runs}); 1.3.2 wrote $\Lambda_\m$ for the free quotient in \S\ref{sec:pg-runs} and for
  the localisation without ``$/\mathrm{tors}$'' in \S\ref{sec:pg-visible}. Each statistic is assigned to its object as the code computes it: $A$
  on the full mod-$\ell$ block, the module route on the full $\mathbb L_\m/\ell^{K}$ (free part separately), saturation and
  \code{hidden\_congruence} on the free quotient of torsion-free blocks only. The lookfor of \S\ref{sec:pg-visible} misread
  \code{torsion\_stable} (a pure-torsion block, not one with a stable free part). \emph{Census:} edition~3.4 replaced by the signed edition~4.0
  (main $+$ b1), with b2 and the L23 section marked as pre-release addenda and the integral census edition~1 listed; every census number
  re-derived from the rows files by the new script \file{docs/latex/tools/pg\_census\_check.py}, which also checks the shape table of
  \S\ref{sec:pg-census34} (now six algebras: the $(1,2,1)$ rows at $n>700$ break ``curvilinear or $\m^{2}=0$''); new Loewy-length-$\ge3$ rows
  at $\ell=59$, $67$, $139$; the saturation defect to $\delta=190$, and $\delta=182$ final at $(700,7)$; the ordinary rows re-joined to the
  records; the DEEP-1 counts as distinct cells ($2158$, $128$) rather than computations ($2575$, $146$). The section label
  \code{sec:pg-census34} is kept (alias \code{sec:pg-census40}). No snippet changed; all twenty outputs regenerated on the slot-25 binary are
  byte-identical (only \file{costs.tex} moved).
\item Version 1.3.2 (slot~24, 2026-10-30). The $(4374,3)$ cell is confirmed: the \code{--sumrules} cross-check resolved and gives \code{max\_tors=7} (report \pfile{docs/implementation/slot24/CELL3_4374_report.md}); the P4 marker is removed.
\item Version 1.3.1 (slot~24, DOCS-FIX, 2026-10-29; referee \pfile{docs/implementation/slot24/DOCS_referee_1029.md}, P1--P9). At
  $\ell=2,3$ the seven-vertex model computes Bredon cohomology for $\Ga$ ({\ttfamily[C,~prop:modeltorsion]}); the L23 data are group
  cohomology on $C_I$ (P1). The Lean status of the $\ell=2,3$ dimensions is stated as in {\ttfamily[C,~rmk:lvdimsltwothree]}: rank formulas for
  every $n$, values cell by cell (P2). $n=0$ excluded from the $H^{1}$ count (P3). The $(4374,3)$ row reads ``main reader only; cross-check
  pending'' until the run's report is committed (P4). \S\ref{sec:pg-torsexp} restructured as definition, theorem, what [C] proves, data, the
  exceptional cells (P5); its history is below. The boundary pair is certified by {\ttfamily[C,~prop:descentpair]} (P6); the third console limit
  $\psi(N)(n+1)\le300$ is named (P7); \S\ref{sec:pg-missing} is a phenomenon$\to$verb table (P8); the level-one Bockstein statement cites
  {\ttfamily[C,~prop:badprimes]} (P9). No snippet or generated output changed; every label kept.

  \emph{History of \S\ref{sec:pg-torsexp}.} Version~1.2.0 stated $e_1=\beta(n,\ell)$ as an empirical observation; it is Callegaro--Cohen--Salvetti's
  theorem. Two earlier forms were wrong, both refuted by the data: ``$e_1=v_2(n)+1$'' fails at every power of $2$ ($10/10$ up to $2048$), and
  ``exponent $1$ at $\ell=3$'' fails at every $n$ with $3\mid n$ ($49$ of the $146$ DEEP-1 cells). The note on the DEEP-1 data stated
  $e_1=v_3(n)+1$, the restriction of $\beta$ to its grid (no $n=2\cdot3^{j}$); both datasets measure the same invariant and agree on the overlap
  $n\in\{600,700\}$. The cells $(1458,3)$, $(2916,3)$, $(4374,3)$ were run in slot~24 (L23S14, catalogue entry PH-01) to check the exception
  beyond DEEP-1. Version~1.2.0 of \S\ref{sec:pg-dihedral3} named only $\ell=59$ as the next onset; $\ell=47$ had been omitted. The table of
  \S\ref{sec:pg-missing} replaced, in slot~22, a list of candidate verbs (\code{level\_n}, \code{nebentypus(N, chi, n, ell)}, \code{lvalue},
  \code{hida\_lambda(ell, class)}, \code{hida\_tower}, \code{kummer\_unit}, \code{forms\_algebra(k, ell, K)}, \code{pow}, a store reader); the
  store reader became the atlas channel, not a verb (owner decision of slot~22).
\item Version 1.3.0 (slot~24, DOCS-PG, 2026-10-29). Reorganised by mathematical phenomenon into six parts (owner direction of 2026-10-29); each
  part shows its snippets, then the run results on the same phenomenon; the former Part~II ``Phenomena from the runs'' is distributed over
  Parts~III--VI, its introduction is \S\ref{sec:pg-runs}. Three new snippets: \file{cohomology\_dimensions} (\S\ref{sec:pg-dims}),
  \file{descent\_torsion\_beta} (\S\ref{sec:pg-descent-torsion}) and \file{descent\_congruence\_primes} (\S\ref{sec:pg-descent-congruence}), the
  integral descent of slot~24. The hand-typed Smith table of \S\ref{sec:pg-l23} is replaced by the generated output; \S\ref{sec:pg-torsexp}
  records the slot-24 cells $(1458,3)$, $(2916,3)$, $(4374,3)$. All outputs regenerated (the seventeen old transcripts unchanged). Every label
  of 1.2.1 is kept.
\item Version 1.2.1 (slot 23, phenomenology package, 2026-10-19, the literature verdicts): Part II cites the sources found by the literature searches (docs/phenomenology/literature\_search\_A.md, \_B.md; bibliography phenomenology\_refs.bib) and corrects the over-statements: the torsion exponent = Frobenius bound is Callegaro--Cohen--Salvetti's Thm 4.5 (KNOWN), (547,486) is Bellaiche--Pollack's exception, (21929,2) e\_2 = 127 is Calegari--Emerton's example, the floor 7 is Rustom's Thm 8.1, (Hi11) rank 3 is Hida control and n = 3460 is the Dickson exception, the dihedral onset list gains ell = 47; cells written (n, ell) throughout. No snippet or generated output changed.
\item Version 1.2.0 (slot 23, phenomenology package, 2026-10-19): Part~II ``Phenomena from the runs'' (\S\ref{sec:pg-runs}--\ref{sec:pg-mazur23}): census 3.4 at a glance, the dihedral Loewy-3 classes, the ordinary Loewy-3 block (547,486), the saturation defect, ell = 2, 3 (L23, DEEP-1), the torsion exponent = Frobenius bound (with the ell = 3 reconciliation), (Hi11) to n = 6910 (second-cluster ACL2O), Mazur's Eisenstein ideal at p = 2, 3; tag [EMP]; no snippet and no generated output changed. Input: docs/phenomenology/phenomena\_catalogue.md.
\item Version 1.1.0 (slot 22, VERB-3, 2026-10-09): one snippet per verb added in slots 21--22 (\code{pow}, \code{kummer\_unit}, \code{hida\_tower}, \code{level\_n}, \code{nebentypus}, \code{forms\_algebra}, \code{hida\_lambda}), \S\ref{sec:pg-missing} rewritten.
\item Version 1.0.0 (hecke-extensions slot 20, stretch package S5, 2026-10-04; owner request of 2026-10-04).
\end{itemize}

\bibliographystyle{alpha}
\bibliography{gl2_refs,phenomenology_refs}

\end{document}

